\documentclass[11pt,a4paper]{article}
\usepackage[a4paper,margin=0.95in]{geometry}
\usepackage[T1]{fontenc}
\usepackage{lmodern,microtype}
\usepackage{amsmath,amssymb,amsthm,mathtools}
\usepackage{booktabs,array,tabularx,longtable,multirow,float}
\usepackage{enumitem,xcolor,url,listings,needspace,multicol}
\usepackage[hidelinks]{hyperref}
\usepackage[nameinlink,noabbrev,capitalise]{cleveref}
\usepackage{fancyhdr}

\definecolor{deepblue}{RGB}{28,69,125}
\hypersetup{colorlinks=true,linkcolor=deepblue,citecolor=deepblue,urlcolor=deepblue,
 pdftitle={Equivalent-Key Recovery and Universal Forgery for QR-UOV: One Regular Direction Determines the Hidden Oil Space},
 pdfauthor={Justin Thaler, Kai Zhang, Scott Kominers, Quang Dao},
 pdfsubject={Equivalent-key recovery, universal forgery, algebraic attacks, and parameter implications for QR-UOV},
 pdfkeywords={QR-UOV,UOV,equivalent-key recovery,universal forgery,structural cryptanalysis,graph direction,Jacobian localization,projective slicing,geometric resolution,exact verification}}

\theoremstyle{definition}
\newtheorem{theorem}{Theorem}[section]
\newtheorem{lemma}[theorem]{Lemma}
\newtheorem{proposition}[theorem]{Proposition}
\newtheorem{corollary}[theorem]{Corollary}
\newtheorem{definition}[theorem]{Definition}
\theoremstyle{remark}
\newtheorem{remark}[theorem]{Remark}

\newcommand{\F}{\mathbb F}
\newcommand{\PP}{\mathbb P}
\newcommand{\cO}{\mathcal O}
\newcommand{\cB}{\mathcal B}
\newcommand{\cX}{\mathcal X}
\newcommand{\cR}{\mathcal R}
\newcommand{\cU}{\mathcal U}
\newcommand{\cV}{\mathcal V}
\newcommand{\cW}{\mathcal W}
\newcommand{\cG}{\mathcal G}
\newcommand{\cE}{\mathcal E}
\newcommand{\Span}{\operatorname{span}}
\newcommand{\rank}{\operatorname{rank}}
\newcommand{\codim}{\operatorname{codim}}
\newcommand{\cdeg}{\operatorname{cdeg}}
\newcommand{\Sym}{\operatorname{Sym}}
\newcommand{\disc}{\mathrm{disc}}
\newcommand{\bad}{\mathrm{bad}}
\newcommand{\oil}{\mathrm{oil}}
\newcommand{\off}{\mathrm{off}}
\newcommand{\spanbad}{\mathrm{span}}
\newcommand{\QRUOV}{\textsf{QR-UOV}}
\newcommand{\Solve}{\operatorname{Solve}}
\newcommand{\Adv}{\operatorname{Adv}}
\newcommand{\trans}{\mathsf T}
\newcommand{\softO}{\widetilde O}
\newcommand{\Mmul}{\mathsf M}
\newcommand{\Gmul}{\mathsf G_{\rm mul}}
\newcommand{\proofstep}[1]{\par\smallskip\noindent\textbf{#1}\enspace}

\setlist[itemize]{leftmargin=*,itemsep=0.10em,topsep=0.20em}
\setlist[enumerate]{leftmargin=*,itemsep=0.12em,topsep=0.20em}
\setlength{\parindent}{1.2em}
\setlength{\parskip}{0.20em}
\linespread{1.04}
\setlength{\emergencystretch}{2em}
\raggedbottom
\pagestyle{fancy}
\fancyhf{}
\fancyfoot[C]{\thepage}
\renewcommand{\headrulewidth}{0pt}

\title{\bfseries Equivalent-Key Recovery and Universal Forgery for \QRUOV{}\\ One Regular Direction Determines the Hidden Oil Space}
\author{Justin Thaler\thanks{Georgetown University and a16z crypto research.}
\and Kai Zhang\thanks{MIT.}
\and Scott Kominers\thanks{Harvard University and a16z crypto research.}
\and Quang Dao\thanks{Carnegie Mellon University and LayerZero Labs.}}
\date{28 July 2026}

\begin{document}
\maketitle

\begin{abstract}
Unbalanced Oil and Vinegar (UOV) hides a linear ``oil'' subspace on which every public quadratic form has no oil--oil term.  Knowing any valid hidden subspace is enough to sign: after fixing the remaining ``vinegar'' variables, the public equations become linear in the oil variables.  \QRUOV{} compresses this structure by working over a small extension field.

We give a direct equivalent-key attack.  Write the hidden oil space as the graph of an unknown matrix $G$.  For any public oil-coordinate direction $c$, the public equations
\[
P_i(-x,c)=0\qquad(1\le i\le m)
\]
have the planted solution $x=Gc$.  If the derivative matrix at that solution has full column rank, then $x$ determines every column of $G$ through public linear systems.  The recovered graph is checked against every public quadratic form, pulled back to the official base-field representation, and tested for an invertible signing system.  An accepted certificate therefore signs every verification target and gives a key-only universal chosen-message forgery; the original secret seed is never needed.

The reduction is both strong and easy to derandomize.  In the public-random expanded-key model, one fixed coordinate direction is regular except with probability about $2^{-62.90}$ at Levels I and III and $2^{-83.86}$ at Level V.  Testing only the first $3,4,4$ coordinate directions, together with an explicit Wronskian rank condenser, gives deterministic lists of just $315,612,1228$ square systems whose rank-exception probabilities are below $2^{-188.69},2^{-251.59},2^{-335.46}$.  For signing, testing $22,31,40$ fixed basis vinegar assignments already gives failure probabilities below $2^{-153.50},2^{-216.30},2^{-279.09}$; a projective-line fallback makes them much smaller.  The basis-first scan uses fewer than $1.008$ tests on average, and its all-basis failure probability is negligible.

We also analyze a faster projective route.  A smallest possible projective slice meets the hidden oil space in one point.  Outside explicit algebraic bad-key sets, a public recombination produces a smooth complete intersection of $V$ quadrics with $2^V$ solutions.  The finite-field Kronecker algorithm of van der Hoeven and Lecerf then gives expected bit complexity
\[
\widetilde O\!\left(2^V2^{(1+\varepsilon)(V-1)}\,3\log_2 127\right)
\]
for every fixed rational $\varepsilon>0$.  We give explicit admissible fields, one-invocation failure bounds, exact output checks, and independent solver cross-checks.

The algebraic recovery and forgery claims are exact.  The key-density statements are exact in the stated public-random and ideal-primitive models.  The normalized asymptotic exponents lie below the three Round-2 category targets, with the most robust margins at Levels III and V.  These are not measured gate counts, and the direct degree-$2^V$ representations require enormous memory.  The paper therefore establishes a structural universal-forgery route and a concrete parameter-review obligation, while leaving implementation-level cost as the main unresolved question.
\end{abstract}

\section{Introduction}\label{sec:intro}

\subsection*{Why recovering an equivalent key is enough}
A multivariate signature scheme publishes a quadratic map
\[
P: x\longmapsto \bigl(P_1(x),\ldots,P_m(x)\bigr)
\]
and keeps a trapdoor for solving $P(x)=y$.  UOV divides the secret coordinates into \emph{vinegar} and \emph{oil} variables.  Its defining feature is that no public equation has a term involving two oil variables.  Once the vinegar variables are fixed, all equations are therefore linear in the oil variables, and signing is one linear-system solve~\cite{KPG1999}.

The public key hides which directions are oil directions.  Cryptanalysis does not have to recover the original seed or the exact serialized private key.  Any hidden oil space that reproduces the public map and admits one invertible oil system is an \emph{equivalent signing key}.  Such a key can be verified directly and can invert every target used by the official verifier.

\QRUOV{} compresses UOV by grouping base-field coordinates into elements of
\[
K=\F_{127^3}.
\]
For the parameter sets studied here, the expanded public key has $m$ quadratic forms in $V+O$ extension-field variables, and the hidden oil space has dimension $O$~\cite{QRUOV2021,QRUOVSpec}.  NIST advanced \QRUOV{} to the third round of the Additional Digital Signatures process in May 2026~\cite{NISTRound3,NISTIR8610}.

\subsection*{The main identity}
After a public normalization, write the hidden oil space as
\begin{equation}\label{eq:intro-graph}
 \cO=\{(-Gc,c):c\in K^O\},
\end{equation}
where $G\in K^{V\times O}$ is unknown.  Choose any public direction $c\in K^O$ and form
\begin{equation}\label{eq:intro-direction}
 q_{i,c}(x):=P_i(-x,c)=0,
 \qquad 1\le i\le m.
\end{equation}
The point $x=Gc$ is automatically a solution.  More importantly, translating by that solution gives the exact formula
\begin{equation}\label{eq:intro-translation}
 q_{i,c}(Gc+z)=z^{\mathsf T}A_i z-2(B_i c)^{\mathsf T}z.
\end{equation}
Thus the derivative matrix at the planted solution is simply the $m\times V$ matrix with rows $(B_i c)^{\mathsf T}$.

If this matrix has full column rank, one recovered root $x=Gc$ reveals the \emph{entire} graph $G$.  For every direction $d\in K^O$, the unknown vector $Gd$ is the unique solution of a public overdetermined linear system.  Taking $d$ to be the $O$ coordinate vectors recovers every column of $G$.  No other root and no global description of the solution set is needed.

The resulting attack has four conceptually simple steps:
\begin{enumerate}
\item choose a public direction $c$ and solve the $m$ equations in \eqref{eq:intro-direction};
\item retain a root whose derivative matrix has full column rank;
\item recover $G$ from $O$ public linear systems; and
\item verify every graph identity, pull the decomposition back to the official base field, and exhibit an invertible oil matrix.
\end{enumerate}
Every accepted output is checked against the unchanged public map.  A polynomial-system solver may fail, omit roots, or return malformed data; those outcomes cause a restart, not a false key.

\subsection*{Small deterministic lists}
The direction system has $m>V$ equations.  A naive square reduction would try all $\binom mV$ row subsets.  We replace that enumeration by a one-parameter Wronskian condenser.  For every full-rank $m\times V$ derivative matrix, a public list of only
\[
V(m-V)+1
\]
linear recombinations contains an invertible $V\times V$ core.  The list lengths are $105,153,307$ at Levels I, III, and V.

In the public-random expanded-key model, the coordinate-direction derivative matrices are independent and uniform.  Consequently, the first $3,4,4$ coordinate directions already give rank-exception probabilities below the corresponding NIST category scales.  Combining these directions with the condenser produces deterministic lists of only
\[
315,
\qquad612,
\qquad1228
\]
square systems.  Scanning every coordinate direction gives still smaller exceptional probabilities, but is not needed for a category-scale statement.

The signing check can also be made shorter.  Certain fixed basis vinegar assignments produce independent uniform $m\times m$ oil matrices over $\F_{127}$.  Testing $22,31,40$ such assignments gives failure probabilities below $2^{-153.50},2^{-216.30},2^{-279.09}$.  Testing all basis assignments makes the failure much smaller, and a final projective-line scan supplies a deterministic fallback.

\subsection*{A faster projective route}
The graph-direction reduction is the cleanest structural result, but its most convenient published solver interface pays for a degree-$V$ localization polynomial.  We therefore also study a faster route.

A projective $V$-plane in the $(V+O-1)$-dimensional ambient space necessarily meets the hidden projective oil space.  For a transverse plane, the intersection is one point.  We prove that, outside three explicit bad-key loci, some public slice and some public recombination of the equations produces a smooth zero-dimensional complete intersection of $V$ quadrics.  It has exactly
\[
D=2^V
\]
geometric solutions.  The planted oil point is one of them, and its common derivative kernel is the full hidden oil space.

The finite-field Kronecker algorithm of van der Hoeven and Lecerf~\cite{vdHL} computes a complete univariate description of such a core in expected bit complexity arbitrarily close to $D^2$.  We state the exact field-size requirement, keep its exponent parameter and field extension linked correctly, bound the failure probability of one finite invocation, and verify the returned representation before using it.  The older straight-line-program bound and the independent locally closed-set theorem of Gim\'enez et al.~\cite{GiustiLecerfSalvy2001,Gimenez2026} provide separate growth-rate checks.

\subsection*{What is proved, and what the cost tables mean}
The paper separates exact algebra from asymptotic cost interpretation.
\begin{center}
\small
\setlength{\tabcolsep}{4.5pt}
\renewcommand{\arraystretch}{1.16}
\begin{tabularx}{0.97\textwidth}{@{}>{\raggedright\arraybackslash\bfseries}p{0.23\textwidth}X@{}}
\toprule
Statement & Meaning \\
\midrule
Recovery from a regular direction & Exact for each fixed expanded key satisfying the stated rank condition.  One suitable root determines a publicly verified equivalent decomposition. \\
Universal forgery & Exact.  Once an invertible oil matrix is exhibited, the recovered key inverts every verification target and signs every chosen message without a signing-oracle query. \\
Public-random key density & Exact in the stated expanded-key distribution.  The ideal-XOF and ideal-cipher appendices transfer these bounds with explicit statistical losses. \\
Solver running time & A published asymptotic bit-complexity theorem, with explicit field and failure accounting. \\
Displayed exponents & Normalized comparisons with the specification's gate-exponent rows.  They are not measurements of a full implementation. \\
Memory & Astronomical for the direct degree-$2^V$ representations.  This is the main obstacle to a practical attack. \\
\bottomrule
\end{tabularx}
\end{center}

The normalized rows are below all three category targets.  The conclusion is strongest at Levels III and V, where several independent solver specializations leave substantial margin.  Level I is more sensitive to hidden constants and representation choices.  Rather than repeating these qualifications throughout the paper, \cref{sec:security} collects them in one place.

\subsection*{Main contributions}
\begin{enumerate}
\item We prove the graph-direction identity \eqref{eq:intro-translation} and show that one regular planted root determines the entire hidden graph through public linear algebra.
\item We give an explicit Wronskian rank condenser, prove its optimal list length, and obtain deterministic category-scale lists of $315,612,1228$ square systems.
\item We give a basis-first signing search using fewer than $1.008$ tests on average, short category-scale deterministic signing lists, and a stronger projective-line fallback.
\item We prove that localizing at one nonzero derivative determinant supplies the regularity and radicality required by a locally closed Kronecker solver.  No assumption is made about singular components outside that open set.
\item We prove a separate projective regularization theorem and explicit public-random bounds for all associated bad-key loci.
\item We give finite-invocation solver failure bounds, verified restart, exact resolution checks, exact graph and pull-back checks, and a one-invocation universal-forgery statement.
\item We distinguish theorem-backed asymptotic costs from diagnostic normalizations, and we quantify the otherwise hidden memory cost.
\end{enumerate}

\paragraph{Terminology.}
A \emph{regular root} is a root whose derivative matrix has full rank.  A \emph{projective slice} is a linear subspace used to reduce the hidden oil space to one point.  A \emph{geometric} or \emph{Kronecker resolution} is a univariate encoding of all roots.  A theorem described as \emph{key by key} applies to every fixed expanded key satisfying its explicit algebraic conditions; it does not assume a key distribution.  The notation $\widetilde O(\cdot)$ suppresses logarithmic factors and constants.

\paragraph{Organization.}
\Cref{sec:model} fixes the expanded-key model and exact signing certificate.  \Cref{sec:recovery} gives graph-direction completion and the short deterministic lists.  \Cref{sec:local-geometry} proves the localization lemma.  \Cref{sec:projective} develops the faster projective route.  \Cref{sec:solver-attack} states the solver interfaces and recovery theorems.  \Cref{sec:security} gives the cost and parameter analysis.  The appendices treat seeded expansion, optional uniqueness, alternative elimination ideas, and reproducibility.

\section{Related work}\label{sec:related}

\paragraph{QR-UOV and generic forgery attacks.}
\QRUOV{} compresses UOV through a quotient-ring representation~\cite{QRUOV2021}.  Its Round-2 specification analyzes direct MQ, reconciliation, intersection, PXL, and rectangular-MinRank attacks~\cite{QRUOVSpec}.  Later work by the \QRUOV{} team consolidates the scheme's security arguments, parameters, and implementations~\cite{FurueEtAl2026}; rectangular-MinRank and truncated-polynomial-ring analyses refine several UOV estimates~\cite{FurueIkematsu,Suzuki,FurueIkematsu2026}.  Generic underdetermined-MQ forgery has also been revisited in \emph{Just Guess} and in generalized variable-partitioning work~\cite{JustGuess,QRUOVJustGuess,Asanuma2026}.  Those attacks seek one signature.  Our endpoint is a publicly verifiable decomposition that signs every target.

\paragraph{Recovering hidden UOV structure.}
Intersection attacks and later UOV cryptanalysis recover information about the hidden oil geometry through structured linear algebra~\cite{Beullens}.  P\'ebereau observed that a suitable oil vector can be enough for polynomial-time expanded-key recovery~\cite{Pebereau}.  Our graph-direction theorem gives a direct completion formula: one regular planted vector determines every column of the hidden graph.  The projective route instead recovers the oil space as a common derivative kernel.

\paragraph{Polynomial-system solving.}
The fastest cost line uses the finite-field Kronecker solver of van der Hoeven and Lecerf~\cite{vdHL}.  Under regularity and geometric radicality of every prefix, it computes a complete resolution in expected time arbitrarily close to quadratic in the number of roots.  We use the older Giusti--Lecerf--Salvy straight-line-program bound as an arithmetic growth check, without importing failure constants that its statement does not provide~\cite{GiustiLecerfSalvy2001}.  Gim\'enez et al. give a probabilistic Kronecker algorithm for locally closed sets over perfect fields~\cite{Gimenez2026}; this is the natural interface for the graph-direction system after localization at a nonzero Jacobian determinant.  The projective density proof also uses Benoist's complete-intersection discriminant degrees~\cite{Benoist}.

\paragraph{Low-space and alternative elimination ideas.}
\Cref{app:alternatives} records exact norm-gradient and resultant-gradient formulas that extract an isolated root from a quotient representation, together with audits of overdetermined Macaulay solving, underslicing, degree smoothing, tracer-based F4, resultants, and Chow forms~\cite{BaurStrassen,DAndreaDickenstein,JeronimoEtAl,NeigerEtAl,DegreeSmoothing,F4Tracer}.  These ideas reduce the desired output but do not currently avoid construction of a degree-$2^V$ algebra.

\section{Expanded-key model and exact signing certificates}\label{sec:model}

This section connects the algebra used by the attack to the official verification map.  There are two coordinate systems:
\begin{itemize}
\item the official base-field system over $\F_{127}$, in which signatures are verified; and
\item a smaller extension-field system over $K=\F_{127^3}$, in which the hidden oil space and the attack are easier to describe.
\end{itemize}
The specification gives an exact, invertible pull-back between them.  We first recover and verify the hidden structure over $K$, then pull it back and check the resulting base-field factorization against the unchanged public key.  One invertible oil system then certifies that the recovered decomposition can sign every target.

The Round-2 $q=127$ parameter sets use quotient-ring degree $\ell=3$.  The specification has $v=\ell V$ base-field vinegar variables and $m=\ell O$ base-field oil variables and output coordinates.  After quotient-ring expansion, the attacker works with the same $m$ outputs as quadratic forms in only $N=V+O$ variables over
\[
K=\F_{q^\ell}=\F_{127^3},\qquad Q=|K|=127^3=2{,}048{,}383,
\qquad \log_2 Q=20.966054\ldots.
\]
Thus $m$ denotes both the number of public output coordinates and the base-field oil dimension, whereas $O=m/\ell$ is the oil dimension over $K$.

Let $\Phi_f:K\to\F_q^{\ell\times\ell}$ be the regular representation from the specification, and let $W_f\in\F_q^{\ell\times\ell}$ be its invertible symmetrizer.  Write $\phi$ for the blockwise inverse of $\Phi_f$, and write $W_f^{(a)}$ for the block-diagonal matrix with $a$ copies of $W_f$.

The attack uses the expanded extension-field tuple
\[
P=(p_1,\ldots,p_m),\qquad p_i(x)=x^{\mathsf T}P_ix,
\qquad x\in K^{N},
\]
where the symmetric matrices $P_i$ share a hidden totally isotropic $O$-space $\cO$:
\begin{equation}\label{eq:isotropy}
u^{\mathsf T}P_iw=0\qquad(u,w\in\cO,\ 1\le i\le m).
\end{equation}
The Round-2 sets analyzed in this paper are
\begin{center}
\begin{tabular}{@{}lrrrrrr@{}}
\toprule
Level & $q$ & $\ell$ & $v$ & $m$ & $V=v/\ell$ & $O=m/\ell$\\
\midrule
I&127&3&156&54&52&18\\
III&127&3&228&78&76&26\\
V&127&3&306&105&102&35\\
\bottomrule
\end{tabular}
\end{center}

\begin{center}
\small
\setlength{\tabcolsep}{5pt}
\begin{tabularx}{0.94\textwidth}{@{}>{\bfseries}p{0.16\textwidth}X@{}}
\toprule
Symbol & Meaning \\
\midrule
$N=V+O$ & Number of extension-field variables. \\
$e=m-V$ & Number of public equations not used in the square $V$-equation core. \\
$D=2^V$ & B\'ezout degree of a smooth projective core. \\
$L_s=\F_{Q^s}$ & Solver working field; the same field supplies slicing, recombination, and rational-root extraction. \\
\bottomrule
\end{tabularx}
\end{center}

All concrete theorems below concern these three sets, for which $m<q$.  The fastest projective route has cost $\softO(2^V2^{(1+\varepsilon)(V-1)}3\log_2 127)$ for every fixed rational $\varepsilon>0$.  The graph-direction route instead uses the explicit locally closed finite-field theorem of Gim\'enez et al.; it is slower but needs only one local rank condition.  Both public-random analyses crucially allow arbitrary correlation between the secret graph and the symmetric public blocks.  \Cref{app:seeded-transfer} proves an exact transfer for the official seeded construction in the SHAKE ideal-XOF and AES ideal-cipher models and isolates the remaining fixed-primitive gap.

\subsection{Public-random expansion model and relation to seeded expansion}

\begin{definition}[Public-random expansion with an arbitrarily correlated secret graph]\label{def:ideal}
First sample
\[
(A_1,\ldots,A_m)\leftarrow\Sym_V(K)^m
\]
from the product-uniform distribution.  Conditional on the resulting tuple $A=(A_i)_i$, sample the secret graph $G\in K^{V\times O}$ according to an arbitrary conditional distribution.  Conditional on $(A,G)$, sample
\[
(E_1,\ldots,E_m)\leftarrow(K^{V\times O})^m
\]
from the product-uniform distribution, set each lower-right public block by the key-generation identity, and expose the resulting expanded public tuple.  Values are sampled once and reused consistently.  No independence between $G$ and $A$ is assumed.
\end{definition}

\begin{proposition}[Exact central-block factorization under graph correlation]\label{prop:ideal-factorization}
In the block notation of \eqref{eq:central}--\eqref{eq:graphid}, define
\[
B_i:=E_i-A_iG.
\]
Conditional on every fixed pair $(A,G)$, the tuple $(B_1,\ldots,B_m)$ is product-uniform and its conditional law does not depend on $(A,G)$.  Consequently, $B$ is independent of $(A,G)$, and the pairs $(A_i,B_i)_{1\le i\le m}$ are independent and uniform on $\Sym_V(K)\times K^{V\times O}$.  The graph $G$ may nevertheless remain arbitrarily correlated with $A$.  The lower-right public blocks $C_i$ are determined by \eqref{eq:graphid}.
\end{proposition}

\begin{proof}
Fix $(A,G)=(a,g)$.  For each $i$, the map
\[
E_i\longmapsto B_i=E_i-a_ig
\]
is a translation of the additive space $K^{V\times O}$ and hence a bijection.  It sends the product-uniform law of $E$ to the product-uniform law of $B$, independently of $(a,g)$.  Thus $B$ is independent of $(A,G)$.  Since $A$ itself is product-uniform, $(A_i,B_i)_i$ has the claimed product law.  The formula for $C_i$ is the key-generation block identity \eqref{eq:graphid}.
\end{proof}

Every distributional calculation used to bound bad keys depends on the central pairs $(A_i,B_i)$ and on fresh public attack coins, not on an independent graph sample.  A secret coordinate change depending on $G$ maps a uniform public slice to a uniform central slice for every fixed key.  Thus arbitrary $G$--$A$ correlation does not change the public-random bad-key bounds.  \Cref{app:seeded-transfer} shows that the official seeded construction instantiates this model, up to explicit statistical losses, in the SHAKE ideal-XOF and AES ideal-cipher models.  The key-by-key theorem applies to fixed-primitive keys as well; what remains unproved there is a high-probability statement about fixed-function key generation.

\subsection{Central and public graph coordinates}

In secret central coordinates $K^V\oplus K^O$,
\begin{equation}\label{eq:central}
F_i=\begin{pmatrix}A_i&B_i\\B_i^{\trans}&0\end{pmatrix},
\qquad F_i(v,o)=v^{\trans}A_iv+2v^{\trans}B_io,
\end{equation}
with $A_i\in\Sym_V(K)$ and $B_i\in K^{V\times O}$ independent and uniform.  The planted oil space is $\cO_0=\{0\}\oplus K^O$.

In normalized public coordinates the hidden space is a graph
\begin{equation}\label{eq:graph}
\cO=\left\{\binom{-Gc}{c}:c\in K^O\right\}.
\end{equation}
Writing
\[
P_i=\begin{pmatrix}A_i&E_i\\E_i^{\trans}&C_i\end{pmatrix},
\]
total isotropy is equivalent to
\begin{equation}\label{eq:graphid}
C_i=G^{\trans}E_i+E_i^{\trans}G-G^{\trans}A_iG.
\end{equation}

\begin{proposition}[Verified graph gives an exact expanded UOV decomposition]\label{prop:verified-key}
Let $G'\in K^{V\times O}$ satisfy every identity \eqref{eq:graphid}, and put
\[
B_i'=E_i-A_iG',\qquad
F_i'=\begin{pmatrix}A_i&B_i'\\ {B_i'}^{\trans}&0\end{pmatrix},\qquad
S_{G'}=\begin{pmatrix}I_V&G'\\0&I_O\end{pmatrix}.
\]
Then
\[
P_i=S_{G'}^{\trans}F_i'S_{G'}
\qquad(1\le i\le m).
\]
Consequently, the specification's inverse pull-back gives an expanded base-field UOV decomposition of the unchanged verification map.
\end{proposition}

\begin{proof}
Since
\[
S_{G'}^{-1}=\begin{pmatrix}I_V&-G'\\0&I_O\end{pmatrix},
\]
a direct block multiplication gives
\[
S_{G'}^{-\trans}P_iS_{G'}^{-1}
=
\begin{pmatrix}
A_i&E_i-A_iG'\\
E_i^{\trans}-{G'}^{\trans}A_i&
C_i-{G'}^{\trans}E_i-E_i^{\trans}G'+{G'}^{\trans}A_iG'
\end{pmatrix}.
\]
The lower-right block is zero by \eqref{eq:graphid}; the remaining blocks are exactly $F_i'$.  Rearranging proves the displayed factorization.

To make the interface with the specification explicit, let hats denote the corresponding base-field matrices.  The forward pull-back of Section~4.1 of the specification has the form
\[
P_i=\phi\!\left((W_f^{(N)})^{-1}\widehat P_i\right),\qquad
F_i'=\phi\!\left((W_f^{(N)})^{-1}\widehat F_i'\right),\qquad
S_{G'}=\phi(\widehat S_{G'}).
\]
Hence the inverse pull-back sets
\[
\widehat P_i=W_f^{(N)}\phi^{-1}(P_i),\qquad
\widehat F_i'=W_f^{(N)}\phi^{-1}(F_i'),\qquad
\widehat S_{G'}=\phi^{-1}(S_{G'}).
\]
The pull-back identity in the specification converts the verified extension-field factorization into
$\widehat P_i=\widehat S_{G'}^{\trans}\widehat F_i'\widehat S_{G'}$ for the unchanged base-field verification map~\cite[Sec.~4.1]{QRUOVSpec}.  These matrix identities publicly certify an exact expanded UOV decomposition.  Signing uses this base-field pull-back, not the extension-field tuple with only $O=m/3$ oil variables.  The separate test in \cref{prop:signing-witness} certifies that the decomposition can invert every target.
\end{proof}

For an expanded base-field decomposition $\widehat P_i=\widehat S^{\mathsf T}\widehat F_i\widehat S$, write the central variables as $(a,o)\in\F_q^v\oplus\F_q^m$.  Because the oil--oil block of every $\widehat F_i$ is zero, fixing $a$ leaves a linear map in the oil variables.  Denote its matrix by
\begin{equation}\label{eq:signing-matrix}
\mathsf L_{\widehat F}(a)\in\F_q^{m\times m},
\qquad
\bigl(\mathsf L_{\widehat F}(a)o\bigr)_i
:=\widehat F_i(a,o)-\widehat F_i(a,0).
\end{equation}

\begin{proposition}[Public signing witness]\label{prop:signing-witness}
If $\det\mathsf L_{\widehat F}(a)\ne0$, then $(\widehat F,\widehat S,a)$ is a publicly checkable signing certificate.  For every target $y\in\F_q^m$, one computes
\[
o=\mathsf L_{\widehat F}(a)^{-1}\bigl(y-\widehat F(a,0)\bigr),
\qquad
x=\widehat S^{-1}(a,o),
\]
and obtains $\widehat P(x)=y$.
\end{proposition}

\begin{proof}
Direct substitution gives $\widehat F(a,o)=y$, followed by $\widehat P(x)=\widehat F(\widehat Sx)=y$.  These signatures verify under the unchanged public map; matching the official signing distribution is neither needed nor claimed.
\end{proof}

\begin{corollary}[Key-only universal chosen-message forgery]\label{cor:universal-forgery}
Given a publicly verified signing certificate from \cref{prop:signing-witness}, an attacker can forge an accepted \QRUOV{} signature on every chosen message without querying a signing oracle.  Concretely, for a message $M$, choose any allowed salt $r$, compute the official target
\[
 \mu=\operatorname{SHAKE256}(\mathsf{seed}_{\rm pk}\mathbin\|\operatorname{BytesToBits}(M),512),
 \qquad y=\operatorname{Hash}(\mu,r),
\]
invert $y$ using \cref{prop:signing-witness}, and output $(r,x)$.  The official verifier accepts because $\widehat P(x)=\operatorname{Hash}(\mu,r)$.
\end{corollary}

\begin{proof}
The specification's verifier recomputes $\mu$ and $\operatorname{Hash}(\mu,r)$ and accepts exactly when the public map evaluated at the signature vector equals that target~\cite[Secs.~4.2--4.3]{QRUOVSpec}.  \Cref{prop:signing-witness} inverts every target under the unchanged public map.  No distributional claim about the forged signature is needed for verification, and the secret seed is never recovered.
\end{proof}

Fix a nonzero $\xi\in\F_q^\ell$.  For each quotient-ring vinegar block $1\le j\le V$, let $a_j$ place $\xi$ in block $j$ and zero elsewhere, and set
\[
M_j:=\mathsf L_{\widehat F}(a_j)\in\F_q^{m\times m}.
\]
For $0\le r\le m$, write
\begin{equation}\label{eq:rank-probability}
\varrho_r:=\Pr_{M\leftarrow\F_q^{m\times m}}[\rank M=r]
=q^{-m^2}\prod_{i=0}^{r-1}\frac{(q^m-q^i)^2}{q^r-q^i}.
\end{equation}
In particular,
\begin{equation}\label{eq:square-invertible}
 \pi_m:=\varrho_m=\prod_{k=1}^m(1-q^{-k}),
 \qquad \sigma_m:=1-\pi_m.
\end{equation}
For the stronger fallback, define
\begin{equation}\label{eq:pair-failure}
\chi_m:=q^{-(m-1)}+(1-q^{-(m-1)})q^{-1},
\qquad
\zeta_{\rm pair}:=
\varrho_{m-1}q^{-1}\chi_m
+(1-\varrho_m-\varrho_{m-1})(1-\varrho_m).
\end{equation}

\begin{proposition}[Basis-first deterministic signing search]\label{prop:signing-search}
For the planted decomposition in the public-random expansion model, the matrices
$M_1,\ldots,M_V$ are independent and uniform in $\F_q^{m\times m}$.
Consequently:
\begin{enumerate}
\item Testing the basis assignments $a_1,a_2,\ldots$ in order uses fewer than
\[
 \pi_m^{-1}<1.008
\]
basis tests on average, stopping at the first success or after all $V$ tests.  It fails to find an invertible matrix exactly when all $V$ basis matrices are singular, an event of probability
\begin{equation}\label{eq:basis-sign-failure}
 \eta_{\rm basis}:=\sigma_m^V,
\end{equation}
whose base-two logarithms at Levels I, III, and V are
\[
-362.823,\qquad-530.280,\qquad-711.692.
\]

\item Much shorter fixed lists already suffice at the category scale.  Testing the first
\[
 b_{\rm sign}=22,\qquad31,\qquad40
\]
basis assignments at Levels I, III, and V fails with probabilities $\sigma_m^{b_{\rm sign}}$, whose base-two logarithms are
\begin{equation}\label{eq:compact-sign-logs}
 -153.502,\qquad-216.298,\qquad-279.095.
\end{equation}

\item Pair the basis matrices as $(M_1,M_2),(M_3,M_4),\ldots$.  Let $\cE_{\rm sign}$ be the event that every designated pair span contains only singular matrices.  Then
\begin{equation}\label{eq:eta-sign}
 \Pr[\cE_{\rm sign}]
 \le \eta_{\rm sign}:=\zeta_{\rm pair}^{\lfloor V/2\rfloor},
\end{equation}
with base-two logarithms below
\[
-544.820,\qquad-796.276,\qquad-1068.687.
\]
Outside $\cE_{\rm sign}$, a deterministic fallback succeeds after at most
\begin{equation}\label{eq:signing-scan-count}
 V+\lfloor V/2\rfloor(q-1)
\end{equation}
oil-matrix tests: first test every $M_j$; if necessary, for each pair $(A,B)$ test $A+tB$ for all $t\in\F_q^\times$.  The basis tests already cover the two remaining projective points $A$ and $B$.
\end{enumerate}

For any fixed exact candidate, define its oil-matrix determinant polynomial
\[
\Delta(c_1,\ldots,c_V):=\det\!\left(\sum_{j=1}^V c_jM_j\right).
\]
Whenever $\Delta$ is nonzero, a uniform scalar-block vinegar assignment is invertible with probability at least $1-m/q$.  The corresponding expected numbers of assignments are at most
\[
\frac{q}{q-m}=1.740,\qquad2.592,\qquad5.773
\]
at Levels I, III, and V.  This last statement is key by key.  For a recovered candidate not known to be planted, every successful determinant test is still an exact public signing certificate.
\end{proposition}

\begin{proof}
Write $b_{i,j,k}\in K=\F_{q^\ell}$ for entry $(j,k)$ of the planted extension-field block $B_i$.  Under the public-random experiment these elements are independent and uniform.  The $(i,k)$ block of the oil coefficient row produced by $a_j$ is
\[
2\xi^{\mathsf T}W_f\Phi_f(b_{i,j,k})\in\F_q^{1\times\ell}.
\]
The factor $2$ is invertible.  The $\F_q$-linear map
\[
K\longrightarrow\F_q^{1\times\ell},
\qquad b\longmapsto \xi^{\mathsf T}W_f\Phi_f(b)
\]
is injective: if $b\ne0$, then $\Phi_f(b)$ is invertible, while $\xi^{\mathsf T}W_f\ne0$.  Equal dimensions make it an isomorphism.  Thus each $M_j$ is uniform, and distinct $j$ use disjoint entries of the independent blocks $B_i$; hence the $M_j$ are independent.

A basis test succeeds with probability $\pi_m$.  Truncating a geometric random variable at $V$ gives expected test count
\[
 \sum_{j=0}^{V-1}\sigma_m^j
 =\frac{1-\sigma_m^V}{\pi_m}<\pi_m^{-1}.
\]
Independence gives $\sigma_m^b$ for every fixed prefix of length $b$, yielding all displayed basis-list probabilities.

It remains to bound one pair.  Let $A,B$ be independent uniform $m\times m$ matrices, and let $E_{A,B}$ be the event that their span contains no invertible matrix.  If $A$ is invertible, this event is impossible.  If $\rank A=m-1$, invertible left and right changes of basis reduce $A$ to $\operatorname{diag}(I_{m-1},0)$ without changing the distribution of $B$.  Write
\[
B=\begin{pmatrix}C&b\\ c^{\mathsf T}&d\end{pmatrix}.
\]
On $E_{A,B}$, the polynomial $\det(A+tB)$ vanishes for every $t\in\F_q$.  Its degree is at most $m<q$, so it is the zero polynomial.  Its linear coefficient is $d$, and, after $d=0$, its quadratic coefficient is $-c^{\mathsf T}b$.  These conditions have probability at most $q^{-1}\chi_m$.  If $\rank A\le m-2$, then $E_{A,B}$ implies that $B$ is singular, an event of probability $1-\varrho_m$.  Conditioning on the rank of $A$ gives $\Pr[E_{A,B}]\le\zeta_{\rm pair}$.  The designated pairs are independent, proving \eqref{eq:eta-sign}.

The basis matrices and $A+tB$ for $t\in\F_q^\times$ enumerate every projective point on each designated pair line, so the deterministic scan succeeds whenever one pair span contains an invertible matrix.  Finally, $\Delta$ has total degree $m$.  If it is nonzero, Schwartz--Zippel gives success probability at least $1-m/q$ for a uniform coefficient vector.
\end{proof}

\section{Graph-direction and derivative-kernel recovery}\label{sec:recovery}

The attack has two local endpoints.  The graph-direction endpoint is the stronger structural statement: one affine root determines the complete hidden graph by public linear algebra.  The derivative-kernel endpoint is retained for the faster projective route.  Both endpoints are followed by the same exact graph, pull-back, factorization, and signing checks.

\subsection{A public graph direction produces a planted root}

Fix a direction $c\in K^O$.  Substitute the public oil coordinate $c$ and the negated vinegar coordinate $-x$ into the public form $P_i$:
\begin{equation}\label{eq:direction-system}
 q_{i,c}(x):=P_i(-x,c)
 =x^{\mathsf T}A_ix-2x^{\mathsf T}E_ic+c^{\mathsf T}C_ic,
 \qquad x\in K^V.
\end{equation}
The system $q_{1,c}=\cdots=q_{m,c}=0$ is public and has only $V$ variables.  Define the public $m\times V$ matrix
\begin{equation}\label{eq:direction-matrix}
 \mathcal D_c(x):=
 \begin{pmatrix}
 (E_1c-A_1x)^{\mathsf T}\\[-1mm]
 \vdots\\[-1mm]
 (E_mc-A_mx)^{\mathsf T}
 \end{pmatrix}.
\end{equation}
Recall $B_i=E_i-A_iG$, and write $\mathcal B(c)$ for the matrix whose $i$th row is $(B_ic)^{\mathsf T}$.

\begin{proposition}[Graph-direction identity]\label{prop:graph-direction-identity}
For every $c\in K^O$, the point
\[
 x_c:=Gc
\]
is a common zero of the public system \eqref{eq:direction-system}.  More precisely, for every $z\in K^V$,
\begin{equation}\label{eq:direction-translate}
 q_{i,c}(Gc+z)=z^{\mathsf T}A_iz-2(B_ic)^{\mathsf T}z.
\end{equation}
Consequently,
\[
 \mathcal D_c(x_c)=\mathcal B(c),
 \qquad
 J_q(x_c)=-2\mathcal B(c).
\]
\end{proposition}

\begin{proof}
Expand \eqref{eq:direction-system} at $x=Gc+z$.  Its constant term is
\[
 c^{\mathsf T}(G^{\mathsf T}A_iG-G^{\mathsf T}E_i-E_i^{\mathsf T}G+C_i)c=0
\]
by \eqref{eq:graphid}.  Its linear term is
$2z^{\mathsf T}(A_iG-E_i)c=-2(B_ic)^{\mathsf T}z$, proving \eqref{eq:direction-translate}.  The two matrix identities follow by differentiation and by $E_i-A_iG=B_i$.
\end{proof}

\begin{proposition}[One regular direction determines the entire graph]\label{prop:graph-direction-completion}
Let $c\in K^O$, let $x$ be a common zero of all the equations $q_{i,c}$, and suppose $\rank\mathcal D_c(x)=V$.  For $d\in K^O$, define
\begin{equation}\label{eq:direction-rhs}
 b_{c,d}(x)_i:=c^{\mathsf T}C_id-x^{\mathsf T}E_id
 \qquad(1\le i\le m).
\end{equation}
Then the overdetermined public system
\begin{equation}\label{eq:direction-completion-system}
 \mathcal D_c(x)y=b_{c,d}(x)
\end{equation}
has at most one solution $y\in K^V$.  If $x=Gc$, its unique solution is $y=Gd$.  Thus solving \eqref{eq:direction-completion-system} for the $O$ coordinate vectors $d=e_j$ recovers $G$ exactly.  Requiring every identity \eqref{eq:graphid} then gives a public polynomial-time certificate of a valid expanded UOV decomposition.
\end{proposition}

\begin{proof}
Full column rank gives uniqueness.  For $x=Gc$, multiply \eqref{eq:graphid} on the left by $c^{\mathsf T}$ and on the right by $d$.  Rearrangement yields
\[
 (E_ic-A_iGc)^{\mathsf T}Gd
 =c^{\mathsf T}C_id-(Gc)^{\mathsf T}E_id,
\]
which is exactly row $i$ of \eqref{eq:direction-completion-system}.  Hence $Gd$ is the unique solution.  The final identities are precisely the hypotheses of \cref{prop:verified-key}.
\end{proof}

\begin{proposition}[Several partial-rank directions can be combined]\label{prop:multi-direction-completion}
Let $c_1,\ldots,c_r\in K^O$ and suppose that roots $x_j=Gc_j$ of the corresponding public direction systems have been found.  Stack their public completion matrices and right-hand sides:
\[
 \mathcal D_{\boldsymbol c}(\boldsymbol x)
 :=\begin{pmatrix}\mathcal D_{c_1}(x_1)\\ \vdots\\ \mathcal D_{c_r}(x_r)\end{pmatrix},
 \qquad
 b_{\boldsymbol c,d}(\boldsymbol x)
 :=\begin{pmatrix}b_{c_1,d}(x_1)\\ \vdots\\ b_{c_r,d}(x_r)\end{pmatrix}.
\]
If the stacked matrix has column rank $V$, then for every $d\in K^O$ the public system
\[
 \mathcal D_{\boldsymbol c}(\boldsymbol x)y=b_{\boldsymbol c,d}(\boldsymbol x)
\]
has the unique solution $y=Gd$.  Thus no individual direction needs full rank: enough complementary derivative information across several recovered direction roots also determines the whole graph.
\end{proposition}

\begin{proof}
For each $j$, the proof of \cref{prop:graph-direction-completion} shows that $Gd$ satisfies the $j$th block of equations.  It therefore satisfies the stacked system.  Full column rank gives uniqueness.
\end{proof}

\begin{remark}[A compact attack certificate]\label{rem:graph-certificate}
A successful graph-direction branch can be certified by the tuple $(c,x,G)$ together with one full-rank minor of $\mathcal D_c(x)$.  A verifier checks the $m$ equations $q_{i,c}(x)=0$, the $O$ linear systems \eqref{eq:direction-completion-system}, all $m$ graph identities, the inverse pull-back, and the signing witness.  No claim about the other roots or about completeness of the solver output enters this certificate.
\end{remark}

\subsection{How often a direction is regular, and how to choose a square core}

Put
\begin{equation}\label{eq:rho-mv}
 \rho_{m,V}(Q):=\prod_{a=0}^{V-1}(1-Q^{a-m}),
 \qquad
 \tau_{m,V}(Q):=1-\rho_{m,V}(Q).
\end{equation}
Thus $\rho_{m,V}(Q)$ is the full-column-rank probability of a uniform $m\times V$ matrix over $K$.

The direction system has $m$ equations but only $V$ variables.  The obvious deterministic reduction to a square system tries all $\binom mV$ row subsets.  The next lemma replaces that list by an optimal-size family of public linear combinations.

Index the $m$ equations by $r=0,\ldots,m-1$.  For $0\le a<V$, write
$(r)_a=r(r-1)\cdots(r-a+1)$, with $(r)_0=1$, and define
\begin{equation}\label{eq:wronskian-condenser}
 \mathsf C(t)_{a+1,r+1}:=(r)_a t^{r-a}
 \qquad(t\in K),
\end{equation}
where the entry is zero when $r<a$.  The first $V$ columns of $\mathsf C(t)$ are triangular with diagonal $0!,1!,\ldots,(V-1)!$.  They are therefore nonsingular because $V-1<127=\operatorname{char}K$.

\begin{lemma}[A one-parameter rank condenser]\label{lem:wronskian-condenser}
Let $B\in K^{m\times V}$ have column rank $V$, where $m\le127$.  Then
\begin{equation}\label{eq:wronskian-polynomial}
 W_B(T):=\det\bigl(\mathsf C(T)B\bigr)
\end{equation}
is a nonzero polynomial of degree at most $V(m-V)$.
\end{lemma}

\begin{proof}
For column $j$ of $B$, form the polynomial
$f_j(T)=\sum_{r=0}^{m-1}B_{r+1,j}T^r$.  The columns of $B$ are independent exactly when the polynomials $f_1,\ldots,f_V$ are independent, and $\mathsf C(T)B$ is their ordinary Wronskian matrix.

Choose another basis $g_1,\ldots,g_V$ of their span with distinct increasing degrees
$0\le d_1<\cdots<d_V<m$.  The leading coefficient of the Wronskian of the $g_j$ is, up to the nonzero product of their leading coefficients,
\[
 \det\bigl((d_j)_a\bigr)_{0\le a<V,\,1\le j\le V}
 =\prod_{i<j}(d_j-d_i).
\]
This is nonzero in $K$ because all $d_j$ are distinct and below $127$.  A change of polynomial basis only multiplies the Wronskian by a nonzero scalar, so $W_B$ is nonzero.  Finally,
\[
 \deg W_B
 \le \sum_{j=1}^V d_j-\frac{V(V-1)}2
 \le V(m-V),
\]
where the last bound uses the $V$ largest possible distinct degrees below $m$.
\end{proof}

\begin{proposition}[Key-by-key sampling and public-random density]\label{prop:graph-direction-density}
For a fixed expanded key, suppose some $V\times V$ minor of $\mathcal B(c)$ is not the zero polynomial in $c$.  A fresh uniform $c\leftarrow K^O$ then satisfies $\rank\mathcal B(c)=V$ with probability at least
\begin{equation}\label{eq:random-direction-success}
 1-\frac{V}{Q}.
\end{equation}
Conditional on this event, a fresh uniform $t\leftarrow K$ makes
$\mathsf C(t)\mathcal B(c)$ invertible with probability at least
\begin{equation}\label{eq:random-recombination-success}
 1-\frac{V(m-V)}{Q}.
\end{equation}

In the public-random model, $\mathcal B(c)$ is a uniform $m\times V$ matrix for every fixed nonzero $c\in K^O$.  In particular, one preselected coordinate direction fails with probability $\tau_{m,V}(Q)$, whose base-two logarithm is
\[
 -62.898161,\qquad -62.898161,\qquad -83.864216
\]
at Levels I, III, and V.  More strongly, the $O$ coordinate-direction matrices
$\mathcal B(e_1),\ldots,\mathcal B(e_O)$ are independent uniform matrices.  Therefore
\begin{equation}\label{eq:all-coordinate-directions-bad}
 \Pr\bigl[\rank\mathcal B(e_j)<V\text{ for every }j\bigr]
 =\tau_{m,V}(Q)^O,
\end{equation}
with base-two logarithms
\[
 -1132.166907,\qquad-1635.352198,\qquad-2935.247544.
\]
\end{proposition}

\begin{proof}
Every maximal minor of $\mathcal B(c)$ is homogeneous of degree $V$ in $c$, so Schwartz--Zippel gives \eqref{eq:random-direction-success}.  Conditional on full rank, \cref{lem:wronskian-condenser} and the univariate root bound give \eqref{eq:random-recombination-success}.

Under \cref{def:ideal}, each $B_i$ is an independent uniform linear map $K^O\to K^V$.  For fixed nonzero $c$, the vector $B_ic$ is uniform in $K^V$, independently across $i$.  For the coordinate directions, the vectors $B_ie_j$ use independent columns of the $B_i$, proving both the fixed-direction law and \eqref{eq:all-coordinate-directions-bad}.  The displayed values are exact evaluations of the rank formulas.
\end{proof}

\begin{corollary}[A short deterministic public list]\label{cor:deterministic-outer-list}
Fix any public list of $V(m-V)+1$ distinct elements of $K$.  Whenever $\mathcal B(c)$ has column rank $V$, at least one $t$ on the list makes the square system
\begin{equation}\label{eq:condensed-direction-system}
 \mathsf C(t)\,(q_{1,c},\ldots,q_{m,c})^{\mathsf T}=0
\end{equation}
have the planted root as a nonsingular solution.

For one fixed coordinate direction, the list sizes at Levels I, III, and V are
\[
 105,\qquad153,\qquad307.
\]
Scanning all coordinate directions uses respectively
\[
 1{,}890,\qquad3{,}978,\qquad10{,}745
\]
square systems, and fails in the public-random model only with probability
$\tau_{m,V}(Q)^O$.  Thus all outer choices can be deterministic; only the polynomial-system solver needs random coins.
\end{corollary}

\begin{proof}
The determinant in \eqref{eq:wronskian-polynomial} has degree at most $V(m-V)$ and cannot vanish at every point of a longer distinct list.  At the planted root, the Jacobian of \eqref{eq:condensed-direction-system} is
$-2\mathsf C(t)\mathcal B(c)$ by \cref{prop:graph-direction-identity}.  The numerical counts use $m-V\in\{2,2,3\}$.
\end{proof}

\begin{corollary}[Category-scale deterministic direction lists]\label{cor:compact-direction-list}
Use the first
\[
 r_{\rm dir}=3,\qquad4,\qquad4
\]
coordinate directions at Levels I, III, and V, and for each direction use the condenser list from \cref{cor:deterministic-outer-list}.  The resulting fully public lists contain
\[
315,\qquad612,\qquad1228
\]
square systems.  In the public-random model, the probability that none has the planted point as a nonsingular root is exactly $\tau_{m,V}(Q)^{r_{\rm dir}}$, with base-two logarithms
\begin{equation}\label{eq:compact-direction-logs}
 -188.694,\qquad-251.593,\qquad-335.457.
\end{equation}
These bounds are already below the Level-I, Level-III, and Level-V category exponents $143,207,272$.  Scanning all $O$ coordinate directions gives the stronger probabilities in \eqref{eq:all-coordinate-directions-bad}.
\end{corollary}

\begin{proof}
The coordinate-direction derivative matrices are independent by \cref{prop:graph-direction-density}.  If one of the first $r_{\rm dir}$ matrices has full column rank, \cref{cor:deterministic-outer-list} supplies a condenser value with nonsingular planted Jacobian.  Multiplying the one-direction failure logarithm by $r_{\rm dir}$ gives \eqref{eq:compact-direction-logs}; multiplying the per-direction list lengths $105,153,307$ gives the displayed system counts.
\end{proof}

\begin{remark}[The list length is geometrically optimal]\label{rem:condenser-optimal}
The Grassmannian of $V$-subspaces of an $m$-dimensional space has dimension $V(m-V)$.  For any fixed full-rank linear map $R:K^m\to K^V$, the subspaces on which $R$ is singular form a hyperplane section in Pl\"ucker coordinates.  Over an algebraic closure, at most $V(m-V)$ such hyperplane sections always have a common point.  Hence no family of at most $V(m-V)$ linear recombinations can work for every full-rank $m\times V$ matrix after all field extensions.  The family in \cref{cor:deterministic-outer-list} meets this natural lower bound exactly.
\end{remark}

\paragraph{Randomized versus deterministic use.}
Sampling one field element $t$ is convenient for the expected-cost theorem.  The short list in \cref{cor:deterministic-outer-list} removes that choice entirely and is far smaller than enumerating all $V$-row subsets.

\subsection{Derivative-kernel recovery for the projective route}

For a public common zero $x$, form
\[
C_x=(P_1x\ \cdots\ P_mx)\in L^{(V+O)\times m}.
\]

\begin{lemma}[Derivative-kernel recovery]\label{lem:kernel}
If $x\in\PP(\cO_L)$, then $\cO_L\subseteq\ker C_x^{\trans}$.  If $\rank C_x=V$, equality holds.
\end{lemma}

\begin{proof}
For $u\in\cO_L$, \eqref{eq:isotropy} gives $u^{\trans}P_ix=0$ for every $i$.  If $\rank C_x=V$, its transpose kernel has dimension $(V+O)-V=O$.
\end{proof}

At the oil point, a smooth restricted core has derivative rank $V$.  Total isotropy bounds the ambient derivative rank by the same value, so \cref{lem:kernel} recovers the complete oil space.  Row reduction into the public graph chart, Frobenius descent to $K$, verification of \eqref{eq:graphid}, and \cref{prop:verified-key} then recover the expanded decomposition.  Exact factorization certifies that decomposition, while \cref{prop:signing-search} supplies an invertible oil matrix whose inverse certifies target inversion.

\paragraph{Optional identification with the planted oil space.}
Neither endpoint needs uniqueness: the algorithm keeps every exactly verified candidate and tests all of them for a signing witness.  \Cref{app:uniqueness} separately shows that a second $K$-rational oil space is extraordinarily unlikely.  Outside that additional event, every accepted oil space is the planted one.

\section{Jacobian localization removes global regularity}\label{sec:local-geometry}

Fix $c\in K^O$ and a public recombination $R\in K^{V\times m}$.  Let
\[
 f=(f_1,\ldots,f_V)^{\mathsf T}:=R(q_{1,c},\ldots,q_{m,c})^{\mathsf T}
\]
and put
\begin{equation}\label{eq:local-jacobian-polynomial}
 h(x):=\det J_f(x).
\end{equation}
The graph-direction condition is local: if $\mathcal B(c)$ has full column rank and $R\mathcal B(c)$ is invertible, then \cref{prop:graph-direction-identity} gives
\begin{equation}\label{eq:h-at-planted}
 h(Gc)=(-2)^V\det(R\mathcal B(c))\ne0.
\end{equation}
No assertion is needed about the system on the complementary hypersurface $h=0$.

\begin{lemma}[Localization gives a reduced regular sequence]\label{lem:jacobian-localization}
Let $k$ be a perfect field, let $f_1,\ldots,f_V\in k[x_1,\ldots,x_V]$, and let $h=\det J_f$.  In the localization
\[
 S_h:=k[x_1,\ldots,x_V]_h,
\]
the sequence $f_1,\ldots,f_V$ is regular.  For every $1\le s\le V$, the ideal $(f_1,\ldots,f_s)S_h$ is radical.  Equivalently, the sequence is reduced and regular on the open set $h\ne0$.
\end{lemma}

\begin{proof}
Let $\mathfrak p$ be a prime of $S_h$ containing the first $s$ equations.  Because $h$ is a unit in $S_h$, the full Jacobian matrix is invertible modulo $\mathfrak p$.  Its first $s$ rows are therefore linearly independent.  The Jacobian criterion over a perfect field shows that the local zero set of $f_1,\ldots,f_s$ is smooth of codimension $s$ at every such prime.  Hence the localized quotient is reduced and every minimal prime has height $s$.  The ring $S_h$ is regular, hence Cohen--Macaulay; an ideal generated by $s$ elements and of height $s$ is a complete intersection.  Thus the generators form a regular sequence, and smoothness gives radicality for every prefix.
\end{proof}

\begin{corollary}[Only the planted simple root is needed]\label{cor:local-hypotheses}
If $\mathcal B(c)$ has full column rank and $R\mathcal B(c)$ is invertible, then $Gc$ belongs to the locally closed set
\[
 \{f_1=\cdots=f_V=0,\ h\ne0\}.
\]
The input tuple $(f_1,\ldots,f_V,h)$ satisfies hypotheses (H1)--(H2) of the locally closed Kronecker algorithm of Gim\'enez et al.  These conclusions are independent of every root and irreducible component contained in $h=0$.
\end{corollary}

\begin{proof}
Equation \eqref{eq:h-at-planted} puts the planted root in the open set.  \Cref{lem:jacobian-localization} is exactly the reduced-regular-sequence condition required after localization.
\end{proof}

\begin{remark}[Why this is a separate geometric route]\label{rem:local-independent}
The local theorem does not use the projective slice, the off-oil incidence, the cumulative-degree bound, the joint discriminant, smoothness of the full projective core, or radicality of components on $h=0$.  A failure in any of those arguments can affect the faster projective line without affecting graph-direction recovery.
\end{remark}

\section{Projective slicing and regular-core geometry}\label{sec:projective}

This section supplies the geometric input to \cref{thm:attack}.  Outside three explicit bad key loci, we show that some slice and output recombination gives a smooth $2^V$-point projective core; we then bound the density of those loci.  The next section chooses a finite working field and turns this existence statement into a verified randomized algorithm.  Secret central coordinates are used only to expose the key distribution.

Let $K=\F_{127^3}$, $Q=|K|$, and
\[
(V,O,m)\in\{(52,18,54),(76,26,78),(102,35,105)\}.
\]
Put $e=m-V\in\{2,2,3\}$ and $N=V+O$.  Work first in secret central coordinates
\[
\cV=K^V\oplus K^O,
\qquad
\cO_0=\{0\}\oplus K^O.
\]
The $m$ central quadrics are
\begin{equation}\label{eq:central-slice}
F_i(v,o)=v^{\mathsf T}A_i v+2v^{\mathsf T}B_i o,
\qquad
A_i\in\Sym_V(K),\ B_i\in K^{V\times O}.
\end{equation}
In the public-random expansion model, the pairs $(A_i,B_i)$ are independent and uniform.  Equivalently, each $F_i$ is uniform in
\[
\cU:=H^0\!\left(\PP(\cV),\mathcal I_{\PP(\cO_0)}(2)\right),
\qquad
D_{\rm form}:=\dim\cU=\frac{V(V+1)}2+VO.
\]
The proof uses central coordinates only to expose the distribution.  The secret change from public to central coordinates lies in $\operatorname{GL}_N(K)$.  Left multiplication by its inverse is a bijection on $L^{N\times(V+1)}$, so a uniform public slice remains uniform in central coordinates.

Smoothness, intersection multiplicity, and derivative-kernel recovery are coordinate invariant.  In particular, the public and central tuples have identically vanishing joint discriminant polynomials simultaneously.  Let
\[
X_F:=V(F_1,\ldots,F_m)\subset\PP(\cV).
\]
It contains the planted oil space $\PP(\cO_0)\simeq\PP^{O-1}$.

\begin{proposition}[Exact marginal law on a fixed transverse slice]\label{prop:planted-slice-law}
Fix, independently of the sampled central forms, a $K$-rational projective $V$-plane $\Lambda$ that meets $\PP(\cO_0)$ transversely in one point $p$.  Under the marginal public-random law of the central pairs $(A_i,B_i)$, the restrictions $F_1|_\Lambda,\ldots,F_m|_\Lambda$ are independent uniform elements of
\[
H^0\!\left(\PP^V,\mathcal I_p(2)\right),
\]
the $K$-vector space of projective quadrics through $p$.
\end{proposition}

\begin{proof}
Choose homogeneous coordinates $[z:t_1:\cdots:t_V]$ on $\Lambda$ with $p=[1:0:\cdots:0]$.  Transversality makes the vinegar projection of $\Lambda/p$ invertible, so after changing the $t$-coordinates its central-coordinate parametrization is
\[
v=t,\qquad o=cz+Lt,
\]
for some nonzero $c\in K^O$ and $L\in K^{O\times V}$.  Hence
\[
F_i|_\Lambda(z,t)
 =2z\,t^{\mathsf T}B_ic
  +t^{\mathsf T}\!\left(A_i+B_iL+L^{\mathsf T}B_i^{\mathsf T}\right)t.
\]
For every prescribed linear coefficient $\ell\in K^V$ and symmetric quadratic coefficient $H\in\Sym_V(K)$, the matrices $B_i$ satisfying $B_ic=\ell$ form an affine space of the same size, and then $A_i=H-B_iL-L^{\mathsf T}B_i^{\mathsf T}$ is forced.  Thus $(A_i,B_i)\mapsto F_i|_\Lambda$ is a surjective linear map with equal-size fibers.  The pairs $(A_i,B_i)$ are marginally independent and uniform by \cref{prop:ideal-factorization}, proving the claim.
\end{proof}

\begin{remark}[Why the statement is marginal]\label{rem:slice-law-marginal}
The model permits the secret graph $G$ to depend arbitrarily on $A$.  Conditioning on $G$ can therefore change the law of $A$, so no conditional-on-$G$ uniformity is asserted.  None is needed: the bad-key events in this section depend only on the marginal central tuple $(A_i,B_i)_i$, while a fresh public slice remains uniform after the fixed secret coordinate change for every individual key.
\end{remark}

\begin{proposition}[Generic isolation by the spare equations]\label{prop:generic-full-isolation}
Fix $p\in\PP^V$ and let $m>V$.  A Zariski-generic $m$-tuple in $H^0(\PP^V,\mathcal I_p(2))^m$ has $p$ as its only common projective zero.  Generically its full Jacobian has rank $V$ at $p$.
\end{proposition}

\begin{proof}
For each $x\ne p$, evaluation at $x$ is a nonzero linear functional on $H^0(\PP^V,\mathcal I_p(2))$.  Requiring all $m$ forms to vanish at $x$ therefore imposes codimension $m$.  The incidence with $x\in\PP^V\setminus\{p\}$ has dimension at most the parameter-space dimension plus $V-m$, strictly smaller than the parameter-space dimension.  By Chevalley's theorem, its image is constructible and its Zariski closure has dimension no larger than the incidence, so that closure is proper.  Hence a generic tuple has no common zero away from $p$.  At $p$, the first derivatives of a uniform quadric through $p$ range over the full cotangent space.  The rank-deficient $m\times V$ derivative matrices form a proper determinantal subvariety.
\end{proof}

The two propositions explain why the overdetermined sliced system generically contains only the planted point.  The main attack nevertheless solves a square core because \cref{prop:generic-full-isolation} is qualitative: it does not supply the finite-field density and solver interface needed for a quantified faster theorem.

Let $L/K$ be the finite field chosen below.  One invocation samples
\[
M\leftarrow L^{N\times(V+1)}.
\]
If $\rank M<V+1$, the trial fails.  Otherwise the attack restricts the public quadrics to the projective $V$-plane $\Lambda_M=\PP(\operatorname{im}M)$:
\[
q_i(t)=F_i(Mt),\qquad t\in\PP^V_L.
\]
It next samples $R\leftarrow L^{V\times m}$.  A rank defect again ends the trial; for full-rank $R$, the square core is
\[
g_a(t)=\sum_{i=1}^m R_{a,i}q_i(t),\qquad 1\le a\le V.
\]
Rank defects are counted as failed draws rather than conditioned away.  This keeps the entries of $(M,R)$ uniform, so the Schwartz--Zippel bound applies directly.  By \cref{lem:theta-ranks}, every rank defect already lies in $\{\Theta_F=0\}$.

In secret affine coordinates, a transverse $\Lambda_M$ is exactly a slice $o=c+Lv$ after translating its unique oil intersection.  The projective parameterization is preferable in the proof because it has no inverses or denominators: every restricted coefficient is simply quadratic in the entries of $M$.

\subsection{Why the full discriminant is enough}

The finite-field Kronecker theorem used below assumes that every affine prefix is a reduced complete intersection.  A smooth full projective core already implies these prefix conditions; no product of prefix discriminants is needed.

\begin{lemma}[A smooth full core gives reduced regular prefixes]\label{lem:full-implies-prefix}
Let $g_1,\ldots,g_V$ be homogeneous quadrics in $\PP^V$ over a field of odd characteristic.  If their common zero scheme is a smooth zero-dimensional complete intersection, then the ordered quadrics form a regular sequence and every homogeneous prefix ideal $(g_1,\ldots,g_s)$ is geometrically radical.  After any projective coordinate change and dehomogenization in an affine chart containing a final root, the affine sequence remains regular and every prefix ideal remains geometrically radical.  These are precisely the solver's hypotheses with localization polynomial equal to $1$.
\end{lemma}

\begin{proof}
Work over an algebraic closure and let $S$ be the homogeneous coordinate ring.  The final ideal has height $V$ and is generated by $V$ elements.  Since $S$ is Cohen--Macaulay, $g_1,\ldots,g_V$ is a regular sequence.  Every prefix $I_s=(g_1,\ldots,g_s)$ is therefore an unmixed complete intersection, and its coordinate ring is Cohen--Macaulay.

Suppose that some prefix were not radical.  A Cohen--Macaulay ring satisfies Serre's condition $(S_1)$, so a nonreduced prefix cannot be generically reduced: some irreducible component occurs in its complete-intersection cycle with multiplicity at least two.  Each remaining quadric is a nonzerodivisor modulo the preceding prefix and hence contains no component.  Until the final cut, every surviving component has positive projective dimension, so every remaining positive-degree hypersurface meets it.  The associativity formula for intersection multiplicities preserves the component's multiplicity factor through these proper cuts.  The final zero-cycle would therefore contain a point with multiplicity at least two, contradicting smoothness of the final scheme.  Thus every homogeneous prefix is geometrically radical.

Localization at a chart equation and dehomogenization preserve regular sequences and reducedness.  The same conclusions hold in every affine chart containing a final root; roots at infinity are simply absent from the affine zero scheme.
\end{proof}

Let $\Delta_{V,V}$ be the reduced discriminant for $V$ quadrics in $\PP^V$.  It vanishes exactly when the common zero scheme is not a smooth complete intersection of codimension $V$.  Because $\operatorname{char}K=127\ne2$, Benoist's arbitrary-characteristic formula~\cite{Benoist} gives the partial degree
\begin{equation}\label{eq:dpartial}
d_V=(V+1)2^{V-1}
\end{equation}
in the coefficients of each equation.  Thus it suffices to make the single polynomial
\begin{equation}\label{eq:theta}
\Theta_F(M,R):=
\Delta_{V,V}\!\left(\sum_iR_{1i}F_i(Mt),\ldots,
\sum_iR_{Vi}F_i(Mt)
\right)
\end{equation}
nonzero.  Each coefficient of a restricted equation has degree at most two in $M$ and one in the corresponding row of $R$.  Since the discriminant has degree $d_V$ in the coefficients of each of the $V$ equations,
\begin{equation}\label{eq:theta-degree}
\deg_{M,R}\Theta_F
\le (2V+V)d_V
=3V(V+1)2^{V-1}.
\end{equation}

\begin{lemma}[Nonvanishing certifies full parameter rank]\label{lem:theta-ranks}
If $\Theta_F(M,R)\ne0$, then $\rank M=V+1$ and $\rank R=V$.
\end{lemma}

\begin{proof}
If $\rank M<V+1$, then $\PP(\ker M)$ is nonempty and every pull-back $F_i(Mt)$, together with all of its first derivatives, vanishes there; the core is singular.  If $\rank R<V$, the output equations are linearly dependent and cannot cut a smooth codimension-$V$ complete intersection.  In either case the full discriminant vanishes.
\end{proof}

\subsection{Geometric nonvanishing of the joint slice-and-output discriminant}

For $x\in X_F$, let
\[
\rho_F(x):=\rank\bigl(dF_1(x),\ldots,dF_m(x)\bigr),
\]
where the differentials are taken on $T_x\PP(\cV)$.

\begin{definition}[Three key-level bad loci]\label{def:bad-loci}
Define:
\begin{align*}
\cB_{\oil}&:=\{\,F:\rho_F(p)<V\text{ for every }p\in\PP(\cO_0)(\overline K)\},\\
\cB_{\off}&:=\{\,F:\exists x\in X_F(\overline K)\setminus\PP(\cO_0)\text{ with }\rho_F(x)<V\},\\
\cB_{\spanbad}&:=\{\,F:\dim\Span(F_1,\ldots,F_m)\le V\}.
\end{align*}
\end{definition}

The next three lemmas separate the proof into linear-section geometry, a second-order separability argument, and generic output mixing.

\Needspace{14\baselineskip}
\begin{lemma}[A transverse finite base slice through a regular oil point]\label{lem:linear-section}
Assume $F\notin\cB_{\oil}\cup\cB_{\off}$.  Choose
$p\in\PP(\cO_0)(\overline K)$ with $\rho_F(p)=V$.  Then a Zariski-open dense set of projective $V$-planes $\Lambda\subset\PP(\cV_{\overline K})$ through $p$ has all of the following properties:
\begin{enumerate}
\item $\Lambda\cap\PP(\cO_0)=\{p\}$;
\item $X_F\cap\Lambda$ is finite;
\item at every $x\in X_F\cap\Lambda$, the restricted differential $T_x\Lambda\to\overline K^m$ has rank $V$.
\end{enumerate}
\end{lemma}

\begin{proof}
\proofstep{Dimension of the base locus.}
First, $\dim X_F\le O-1$.  Otherwise choose a closed off-oil point $x$ on a component of dimension at least $O$.  The Zariski tangent space is the common differential kernel and has dimension at least the local dimension, so
\[
\rho_F(x)=(N-1)-\dim T_xX_F\le V-1,
\]
contrary to $F\notin\cB_{\off}$.  Since $\PP(\cO_0)\subset X_F$, the maximal dimension is exactly $O-1$.

\proofstep{Behavior at the planted point.}
At $p$, every differential annihilates $T_p\PP(\cO_0)$.  Both this tangent space and $\ker dF(p)$ have dimension $O-1$, so they are equal.  Therefore a general $V$-plane through $p$ is transverse at $p$ and intersects the oil space only in $p$.  We henceforth restrict to this nonempty open subset of the Grassmannian; no other oil point can occur in the slice.

\proofstep{Off-oil points not lying on a base line.}
Let $\mathcal G_p$ be the Grassmannian of projective $V$-planes through $p$, of dimension $V(O-1)$.  We now consider only $x\in X_F\setminus\PP(\cO_0)$.  Put $K_x=\ker dF(x)$; the off-oil hypothesis gives $\dim K_x\le O-1$.  For fixed $x$, planes through $p$ and $x$ form a family of dimension $(V-1)(O-1)$.  If the tangent direction of the line $\overline{px}$ is not in $K_x$, the additional condition $T_x\Lambda\cap K_x\ne0$ is a proper Schubert condition and lowers this dimension by at least one.  Allowing $x$ to vary in dimension at most $O-1$ gives total dimension at most $V(O-1)-1$.

\proofstep{Lines through the planted point.}
If the tangent direction of $\overline{px}$ lies in $K_x$, then every quadric $F_i$ vanishes identically on that line: both endpoint values and the polar cross term vanish.  Each such direction contributes a full projective line contained in $X_F$, so passing from points of $X_F$ to line directions lowers dimension by at least one.  The direction variety therefore has dimension at most $\dim X_F-1\le O-2$.  Planes containing one such line again form a family of dimension at most
\[
(O-2)+(V-1)(O-1)=V(O-1)-1.
\]
\proofstep{Conclusion.}
The off-oil rank-defect incidence is locally closed and has dimension at most $V(O-1)-1$ by the two cases above.  Its closure in the projective product $X_F\times\mathcal G_p$ has the same dimension, and the projection of that closure is therefore a proper closed subset of $\mathcal G_p$.  Outside it, every intersection point has restricted derivative rank $V$.  At any such point the Zariski tangent space of $X_F\cap\Lambda$ is $T_x\Lambda\cap\ker dF(x)=0$.  Since local dimension never exceeds tangent-space dimension, no point can lie on a positive-dimensional component.  Hence the intersection is finite.
\end{proof}

\Needspace{11\baselineskip}
\begin{lemma}[A second-order spare equation forces separability]\label{lem:second-order}
Let $f_1,\ldots,f_m$ be quadrics on $\PP^V$ with a common point $p$.  Suppose their differential matrix at $p$ has rank $V$.  If there is a nonzero
\[
h\in\Span(f_1,\ldots,f_m)
\]
whose first jet at $p$ is zero, then the rational map
\[
\phi_f:\PP^V\dashrightarrow\PP^{m-1},
\qquad
x\longmapsto[f_1(x):\cdots:f_m(x)]
\]
has differential rank $V$ on a nonempty open subset.
\end{lemma}

\begin{proof}
Choose projective coordinates with $p=[1:0:\cdots:0]$ and affine coordinates $y=(y_1,\ldots,y_V)$.  After an equation-basis change, the first $V$ forms have expansions
\[
f_j(y)=y_j+q_j(y),
\]
where each $q_j$ is quadratic.  Since $h$ is a quadric with zero value and zero differential at $p$, it is a nonzero homogeneous quadratic $q(y)$ on this chart.

Consider the $(V+1)\times(V+1)$ first-jet minor whose rows are $f_1,\ldots,f_V,h$ and whose columns are value followed by the $V$ first derivatives.  Up to a nonzero sign, its lowest-degree homogeneous term is
\[
q(y)-\nabla q(y)\cdot y=-q(y),
\]
using Euler's identity $\nabla q(y)\cdot y=2q(y)$.  It is therefore nonzero.  Hence the full first-jet rank is $V+1$ generically, which is equivalent to differential rank $V$ for $\phi_f$ away from the base locus.
\end{proof}

\begin{lemma}[Finite regular base locus plus separability]\label{lem:finite-base-mixing}
Let $f_1,\ldots,f_m$ be quadrics on $\PP^V$ over an algebraically closed field.  Assume:
\begin{enumerate}
\item their common base scheme $B=V(f_1,\ldots,f_m)$ is finite and nonempty;
\item the $m\times V$ projective differential matrix has rank $V$ at every point of $B$;
\item $\phi_f$ has differential rank $V$ on a nonempty open subset.
\end{enumerate}
Then a nonempty Zariski-open set of matrices $R\in k^{V\times m}$ makes
\[
\left(\sum_iR_{1i}f_i,\ldots,\sum_iR_{Vi}f_i\right)
\]
a smooth zero-dimensional complete intersection of degree $2^V$.
\end{lemma}

\begin{proof}
Let $\cR=k^{V\times m}$.  For $x\notin B$, write
\[
a(x)=(f_1(x),\ldots,f_m(x))\ne0
\]
and let $J(x)$ be the differential matrix.  A matrix $R$ has $x$ as a root exactly when $Ra(x)=0$, a codimension-$V$ linear condition on $\cR$.

On the open set where the first-jet matrix $[a(x)\mid J(x)]$ has rank $V+1$, the map
\[
a(x)^\perp\longrightarrow T_x^*\PP^V,
\qquad
r\longmapsto rJ(x)
\]
is surjective.  Consequently the additional condition $\det(RJ(x))=0$ has codimension one inside $Ra(x)=0$.  After allowing $x$ to vary in a $V$-dimensional space, this singular-root incidence still has codimension at least one in $\cR$.

The first-jet rank-defect locus away from $B$ has dimension at most $V-1$ by assumption~3.  There the root condition $Ra(x)=0$ alone has codimension $V$, so its total incidence also has dimension at most $\dim\cR-1$.

Finally, at each $x\in B$, the root condition is automatic, but the rank-$V$ derivative assumption makes $\det(RJ(x))=0$ a genuine hypersurface in $\cR$.  Since $B$ is finite, these basepoint hypersurfaces do not fill $\cR$.

Globally, the singular-root incidence is the closed subset of $\PP^V\times\cR$ cut out by $Ra(x)=0$ and the maximal minors of $RJ(x)$.  Its projection is closed because $\PP^V$ is proper, and the preceding stratified dimension bounds show that this projection is proper.  Choose $R$ outside it.  The resulting projective zero scheme is nonempty because it contains $B$.  It cannot have positive dimension: at every root the Jacobian has rank $V$, so the Zariski tangent space is zero-dimensional.  Therefore the zero scheme is smooth and zero-dimensional.  Since it is cut out by $V$ quadrics in $\PP^V$, it is a complete intersection of degree $2^V$.
\end{proof}

\begin{theorem}[Existence of a smooth projective core]\label{thm:joint-nonzero}
If
\[
F\notin\cB_{\oil}\cup\cB_{\off}\cup\cB_{\spanbad},
\]
then the joint polynomial $\Theta_F(M,R)$ in \eqref{eq:theta} is nonzero.
\end{theorem}

\begin{proof}
Choose a regular geometric oil point $p\in\PP(\cO_0)(\overline K)$.  By \cref{lem:linear-section}, a dense open subset of the irreducible Grassmannian of $V$-planes through $p$ has finite base intersection and full restricted derivative rank at every basepoint.

After scalar extension, let
\[
T_{F,\overline K}:\overline K^m\longrightarrow\cU_{\overline K},\qquad
(\lambda_i)_i\longmapsto\sum_i\lambda_iF_i.
\]
The coefficient kernel of the derivative map at $p$ has dimension $m-V=e$ over $\overline K$.  Scalar extension preserves the rank of the coefficient map, so
\[
\dim_{\overline K}\ker T_{F,\overline K}
=m-\dim_K\Span_K(F_1,\ldots,F_m).
\]
This relation space lies inside the derivative kernel.  Since $F\notin\cB_{\spanbad}$, its dimension is strictly smaller than $e$.  We may therefore choose a derivative-kernel vector outside $\ker T_{F,\overline K}$; it produces a nonzero global quadric $H\in\Span_{\overline K}(F_i)$ with zero first jet at $p$.

The condition $H|_\Lambda\not\equiv0$ is open and nonempty: if $H$ vanished on every $V$-plane through $p$, it would vanish on their union, which is all of $\PP(\cV)$.  Intersecting this open set with the one supplied by \cref{lem:linear-section} gives a plane $\Lambda$ satisfying both properties.  On this plane, \cref{lem:second-order} makes the restricted linear system generically separable, and \cref{lem:finite-base-mixing} supplies an output matrix $R$ whose core is a smooth complete intersection of degree $2^V$.  Hence $\Theta_F$ is not the zero polynomial.
\end{proof}

\subsection{Finite-field density of the three bad loci}

\subsubsection{First-jet surjectivity away from oil}

\begin{lemma}[The central-form space spans every off-oil first jet]\label{lem:jet-surj}
For every $x\in\PP(\cV)\setminus\PP(\cO_0)$, the first-jet evaluation map
\[
\cU\longrightarrow K\oplus T_x^*\PP(\cV)
\]
is surjective after scalar extension to $\overline K$.
\end{lemma}

\begin{proof}
Choose a vinegar coordinate $v_j$ that is nonzero at $x$ and normalize to the chart $v_j=1$.  The quadrics
\[
v_j^2,\quad v_jv_k\ (k\ne j),\quad v_jo_a\ (1\le a\le O)
\]
belong to $\cU$ and restrict to the affine functions $1$, the remaining vinegar coordinates, and the oil coordinates.  Their first jets form a basis of the value-and-cotangent space.
\end{proof}

\Needspace{8\baselineskip}
\subsubsection{The off-oil incidence}

The codimension count has a simple structure.  At a fixed off-oil point, the values and first derivatives of each central form can be prescribed independently by \cref{lem:jet-surj}.  Making the point a common zero costs $m$ conditions, and forcing derivative rank at most $V-1$ adds the usual determinantal codimension.  Projecting away the point can recover at most its $N-1$ coordinates.  The next lemma makes this count uniform and controls the projection degree.

Let $\mathcal A=\cU^m$ be the affine key space.  For an affine variety, write $\cdeg$ for the sum of the degrees of its irreducible components.  For each vinegar chart $C_j\simeq\mathbb A^{N-1}$, let $Z_j\subset\mathcal A\times C_j$ be the incidence of pairs $(F,x)$ such that all $F_i(x)$ vanish and the $m\times(N-1)$ derivative matrix has rank at most $V-1$.

\begin{lemma}[Codimension and cumulative degree of the off-oil incidence]\label{lem:off-incidence}
Each $Z_j$ is pure of codimension
\[
m+c_{\det},\qquad c_{\det}=(m-V+1)O=(e+1)O,
\]
and satisfies
\[
\cdeg(Z_j)\le 3^m(2V)^{c_{\det}}.
\]
If $Y_j$ is the Zariski closure of its projection to key space, then
\[
\codim_{\mathcal A}Y_j\ge c_{\off}:=m+c_{\det}-(N-1)=e(O+1)+1,
\qquad
\cdeg(Y_j)\le\cdeg(Z_j).
\]
\end{lemma}

\begin{proof}
\proofstep{Structure and codimension.}
On the chart $v_j=1$, the monomials $v_j^2$, $v_jv_k$ for $k\ne j$, and $v_jo_a$ restrict to $1$, the remaining vinegar coordinates, and the oil coordinates.  They give a polynomial splitting of value-and-first-jet evaluation.  After a triangular polynomial change of coefficient coordinates, $Z_j$ is the product of an affine space, the equations setting $m$ value coordinates to zero, and the determinantal variety of $m\times(N-1)$ matrices of rank at most $V-1$.  Hence $Z_j$ is irreducible and pure of codimension
\[
m+(m-V+1)((N-1)-(V-1))=m+(m-V+1)O.
\]

\proofstep{Cumulative degree.}
In the original coefficient-and-point coordinates, every value equation has total degree at most three.  Every derivative entry has degree at most two, so each $V\times V$ minor has degree at most $2V$.  At a generic rank-$(V-1)$ point, the determinantal variety is smooth of codimension $c_{\det}$, and the differentials of its maximal minors span the conormal space.  Choose $c_{\det}$ generic linear combinations of those minors whose differentials form a basis there.  Together with the $m$ value equations, they cut a complete intersection of codimension $m+c_{\det}$ that contains the irreducible variety $Z_j$ as a component.  Cumulative B\'ezout therefore gives
\[
\cdeg(Z_j)\le 3^m(2V)^{c_{\det}}
\]
by the cumulative-degree bounds of Heintz and Lachaud--Rolland~\cite{Heintz,LachaudRolland}.

\proofstep{Projection to key space.}
Projection can recover at most the $N-1$ chart coordinates, which gives
\[
\codim_{\mathcal A}Y_j\ge m+c_{\det}-(N-1)=e(O+1)+1.
\]
The coordinate projection $\mathcal A\times C_j\to\mathcal A$ is linear.  The standard degree-under-linear-projection inequality gives
\[
\deg \overline{\pi(Z_j)}\le\deg Z_j
\]
for an irreducible affine variety~\cite{Heintz,LachaudRolland}.  Since both $Z_j$ and its image closure $Y_j$ are irreducible, this is exactly $\cdeg(Y_j)\le\cdeg(Z_j)$.
\end{proof}

\begin{proposition}[Explicit off-oil density]\label{prop:off-density}
Under the ideal-expansion distribution,
\begin{equation}\label{eq:etaoff}
\Pr[F\in\cB_{\off}]
\le
\eta_{\off}:=
V\,3^m(2V)^{(e+1)O}
Q^{-\{e(O+1)+1\}}.
\end{equation}
\end{proposition}

\begin{proof}
Every off-oil point lies in one of the $V$ vinegar charts, so $\cB_{\off}\subseteq\bigcup_jY_j$.  The affine finite-field point bound
\[
|Y_j(\F_Q)|\le\cdeg(Y_j)Q^{\dim Y_j}
\]
from Lachaud--Rolland~\cite{LachaudRolland}, followed by the union bound and \cref{lem:off-incidence}, gives the formula after division by $|\mathcal A(\F_Q)|$.
\end{proof}

\subsubsection{The oil and low-span loci}

Let
\[
\tau_{m,V}=1-\prod_{j=0}^{V-1}(1-Q^{j-m})
\]
be the rank-deficiency probability of a uniform $m\times V$ matrix.  If every geometric oil direction is rank deficient, then in particular all $O$ standard oil-coordinate directions are rank deficient.  Their derivative matrices are independent in the public-random expansion model.  Hence
\begin{equation}\label{eq:etaoil}
\Pr[F\in\cB_{\oil}]\le\eta_{\oil}:=\tau_{m,V}^{O}.
\end{equation}

Let $e=m-V$ and let $T_F:K^m\to\cU$ be the coefficient map from the proof of \cref{thm:joint-nonzero}.  The event $F\in\cB_{\spanbad}$ is exactly $\dim\ker T_F\ge e$.  For each fixed $e$-space $W\subset K^m$, the condition $W\subseteq\ker T_F$ imposes $eD_{\rm form}$ independent linear conditions on the uniform map $T_F$.  A union bound over the Grassmannian therefore gives
\begin{equation}\label{eq:etaspan}
\Pr[F\in\cB_{\spanbad}]
\le
\eta_{\rm span}:=
\genfrac{[}{]}{0pt}{}{m}{e}_Q Q^{-eD_{\rm form}}.
\end{equation}
Thus the geometric proof requires only $\dim\Span(F_1,\ldots,F_m)>V$; full linear independence is unnecessary.

Combining \cref{thm:joint-nonzero} with
\eqref{eq:etaoff}, \eqref{eq:etaoil}, and \eqref{eq:etaspan} gives the key theorem.

\Needspace{15\baselineskip}
\begin{theorem}[Almost-all-key joint regularizability]\label{thm:key-density}
In the public-random expansion model,
\begin{equation}\label{eq:keybad-total}
\Pr[\Theta_F\equiv0]
\le\eta_{\off}+\eta_{\oil}+\eta_{\rm span}.
\end{equation}
For the Round-2 parameter sets:
\begin{center}
\begin{tabular}{@{}lrrrr@{}}
\toprule
Level & $c_{\off}$ & $\log_2\eta_{\off}$ & $\log_2\eta_{\oil}$ & $\log_2\Pr[\Theta_F\equiv0]$\\
\midrule
I   & 39  & $-364.564$  & $-1132.167$ & $<-364.563$\\
III & 55  & $-457.920$  & $-1635.352$ & $<-457.919$\\
V   & 109 & $-1038.067$ & $-2935.248$ & $<-1038.066$\\
\bottomrule
\end{tabular}
\end{center}
The low-span term is below $2^{-94850}$ already at Level I and is invisible at the displayed precision.
\end{theorem}

\begin{proof}
By \cref{thm:joint-nonzero}, the event $\Theta_F\equiv0$ is contained in $\cB_{\oil}\cup\cB_{\off}\cup\cB_{\spanbad}$.  Apply \eqref{eq:etaoff}, \eqref{eq:etaoil}, \eqref{eq:etaspan}, and the union bound.  The displayed values are direct evaluations of those formulas, rounded outward at the shown precision.
\end{proof}

\section{Finite-field solver interfaces and verified restart}\label{sec:solver-attack}

We first specialize a finite-field theorem for locally closed sets to the graph-direction system.  This gives a recovery route whose correctness needs only the local Jacobian condition of \cref{sec:local-geometry}.  We then return to the faster projective route and its near-quadratic solver.  In both cases, every accepted key is checked against the unchanged public map.

\subsection{Locally closed finite-field solving}\label{subsec:local-solver}

The following proposition records the exact external interface used below.  It is a specialization of Theorem~12 of Gim\'enez et al.~\cite{Gimenez2026}.

\begin{proposition}[Locally closed finite-field solver]\label{prop:gimenez-theorem}
Let $F_1,\ldots,F_r,H\in\F_{\mathfrak q}[X_1,\ldots,X_n]$ have degree at most $d$.  Suppose that, in the localization $\F_{\mathfrak q}[X]_H$,
\begin{enumerate}
\item $F_1,\ldots,F_r$ form a regular sequence; and
\item $(F_1,\ldots,F_s)$ is radical for every $1\le s\le r$.
\end{enumerate}
Let $V_s$ be the Zariski closure of
$\{F_1=\cdots=F_s=0,\ H\ne0\}$, put $\delta_s=\deg V_s$, and let $\delta=\max_s\delta_s$.  If the input has division-free straight-line-program length $L$, $0<\epsilon<1/4$, and
\begin{equation}\label{eq:gimenez-field-condition}
 \mathfrak q>2\epsilon^{-1}n^2rd\delta^3,
\end{equation}
then a probabilistic algorithm computes a Kronecker representation of a lifting fiber of $V_r$ with probability at least $1-2\epsilon$ and bit complexity
\begin{equation}\label{eq:gimenez-cost}
 \widetilde O\!\left(
   \bigl(n^4+r(L+n^2+r^4)\bigr)d\delta^2\log_2\mathfrak q
 \right).
\end{equation}
When $r=n$, the lifting fiber is the entire zero-dimensional locally closed set.  The solver is Monte Carlo: its paper does not provide an efficient certificate that a returned representation is complete.
\end{proposition}

\begin{proof}
The two displayed algebraic conditions are hypotheses (H1)--(H2) of the cited theorem.  Equation \eqref{eq:gimenez-field-condition}, success probability $1-2\epsilon$, and \eqref{eq:gimenez-cost} are its finite-field statement.  The source explicitly notes that its algorithms do not appear to be Las Vegas because output correctness cannot in general be verified at comparable cost.  For $r=n$, no free projection coordinates remain, so a lifting fiber is the zero-dimensional set itself.  The source treats $\delta>d$ as without loss of generality; if $\delta\le d$, replacing it by the safe upper bound $d+1$ only enlarges the field and cost bounds.
\end{proof}

\begin{corollary}[Graph-local specialization]\label{cor:gimenez-local}
Assume $\rank\mathcal B(c)=V$ and $R\mathcal B(c)$ is invertible.  Apply \cref{prop:gimenez-theorem} to the $V$ affine quadrics $f=Rq_c$ and the localization polynomial $h=\det J_f$.  Then
\[
 n=r=V,\qquad d=V,\qquad \delta\le D:=2^V.
\]
The input has straight-line-program length
\begin{equation}\label{eq:local-slp}
 L_{\rm loc}=O(L_{\rm core}+V^4),
 \qquad
 L_{\rm core}:=M_V+2V(M_V+V+1),
 \quad M_V:=\frac{V(V+1)}2.
\end{equation}
Over any extension $L=\F_{Q^s}$ satisfying
\begin{equation}\label{eq:local-field-condition}
 Q^s>2\epsilon^{-1}V^4 8^V,
\end{equation}
the solver contains the planted root $Gc$ with probability at least $1-2\epsilon$ and has bit complexity
\begin{equation}\label{eq:local-solver-cost}
 \widetilde O\!\left(
 \bigl(V^4+V(L_{\rm loc}+V^2+V^4)\bigr)V4^V\log_2(Q^s)
 \right).
\end{equation}
\end{corollary}

\begin{proof}
\Cref{cor:local-hypotheses} gives the localized regularity and radicality hypotheses and places $Gc$ in the open set.  B\'ezout bounds every intermediate degree by $2^s\le2^V$.  The maximum input degree is $V$ because $h$ has degree $V$.  The dense quadratic core has the shared-monomial circuit in \eqref{eq:local-slp}; Baur--Strassen differentiation followed by a division-free Samuelson--Berkowitz determinant circuit adds $O(V^4)$ operations.  Substitution into \eqref{eq:gimenez-field-condition}--\eqref{eq:gimenez-cost} proves the result.
\end{proof}

\begin{corollary}[Smooth-projective Gim'enez cross-check]\label{cor:gimenez-projective}
On a smooth projective core, choose a hyperplane at infinity and dehomogenize in the complementary affine chart.  Conditional on the hyperplane avoiding all $D=2^V$ projective roots, apply \cref{prop:gimenez-theorem} with $H=1$.  Then
\[
 n=r=V,\qquad d=2,\qquad\delta\le D,
\]
and, over $\F_{Q^s}$ with
\begin{equation}\label{eq:projective-gimenez-field}
 Q^s>4\epsilon^{-1}V^3 8^V,
\end{equation}
the direct displayed complexity is
\begin{equation}\label{eq:projective-gimenez-cost}
 \widetilde O\!\left(
  2\bigl(V^4+V(L_{\rm core}+V^2+V^4)\bigr)4^V\log_2(Q^s)
 \right).
\end{equation}
The conditional solver success probability is at least $1-2\epsilon$.  If the hyperplane is sampled uniformly over $\F_{Q^s}$, the unconditional success probability of the chart-plus-solver step is at least
\begin{equation}\label{eq:projective-gimenez-success}
 1-2\epsilon-\frac{D}{Q^s}.
\end{equation}
\end{corollary}

\begin{proof}
Smoothness of the full projective zero scheme implies regularity and geometric radicality of every prefix by \cref{lem:full-implies-prefix}.  A hyperplane avoiding the finite zero set exists, and dehomogenization in its complementary affine chart preserves those properties.  Substitute $d=2$ and $\delta\le2^V$ into \cref{prop:gimenez-theorem}.  A uniformly random hyperplane contains any fixed projective point with probability at most $1/Q^s$; a union bound over the $D$ roots gives the additional $D/Q^s$ term in \eqref{eq:projective-gimenez-success}.
\end{proof}

\paragraph{Why Monte Carlo output is sufficient here.}
A successful call to \cref{prop:gimenez-theorem} contains the planted nonsingular root.  A failed call may omit it or return unusable data.  The attack parses every returned representation, retains only actual roots of all original equations, completes and verifies the graph, checks the inverse pull-back and public-map factorization, and finally checks an invertible oil system.  If no candidate passes, it retries with fresh coins.  Thus the external solver's lack of a completeness certificate can delay termination but cannot create a false accepted key.

\subsection{Fast projective route: field size, failure probability, and cost}\label{subsec:field}

Put
\[
 D:=2^V,\qquad L_s:=\F_{Q^s},\qquad q_s:=|L_s|.
\]
The finite-field theorem of van der Hoeven and Lecerf has an exponent parameter and a field-size requirement.  They must be chosen together.  We use $\varepsilon>0$ for the desired final exponent and set the theorem's internal parameter to $\eta=\varepsilon/3$.  For $V$ quadrics, its field-size threshold is
\begin{equation}\label{eq:vdhl-field-bound}
 B_V(\eta):=
 \max\!\left\{
 \max\!\left\{2,(1+\eta)(8/\eta+1)V2^{\eta/8}\right\}(2^V-1),
 100\,8^V
 \right\}.
\end{equation}
Define the smallest admissible extension degree and a one-degree safety choice by
\begin{equation}\label{eq:solver-extension-degree}
 s_\varepsilon^0:=\min\{s\ge1:Q^s>B_V(\varepsilon/3)\},
 \qquad
 s_\varepsilon^+:=s_\varepsilon^0+1.
\end{equation}
The first choice minimizes field size.  The second divides all finite-invocation failure bounds by the large factor $Q=2{,}048{,}383$.

For any $s\ge s_\varepsilon^0$, sample
\[
 M\leftarrow L_s^{(V+O)\times(V+1)},
 \qquad
 R\leftarrow L_s^{V\times m}
\]
and form the restricted square core.  For every fixed key with $\Theta_F\not\equiv0$, Schwartz--Zippel gives
\begin{equation}\label{eq:fast-beta}
 \beta_{V,s}:=\Pr[\Theta_F(M,R)=0]
 \le \frac{E_V^{\rm joint}}{Q^s},
 \qquad
 E_V^{\rm joint}:=3V(V+1)2^{V-1}.
\end{equation}
The proof of the finite solver theorem bounds its three internal bad choices by
\begin{equation}\label{eq:fast-alpha}
 \alpha_{V,s}:=\frac{5D^2+26D^3}{Q^s}.
\end{equation}
These are bounds for one finite invocation.  A denominator failure, degree mismatch, or failed final check returns \textsf{fail}; the algorithm never runs an unproved internal restart loop on an irregular outer draw.

For the primary choice $\varepsilon=0.3$, the field threshold is governed by $100\,8^V$.  The minimum and safety choices are:
\begin{center}
\small
\begin{tabular}{@{}lrrrr@{}}
\toprule
Level & $s_{0.3}^0$ & success at $s_{0.3}^0$ & $s_{0.3}^+$ & success at $s_{0.3}^+$ \\
\midrule
I   &8  &$\ge0.992337625602$ &9  &$\ge0.999999996259$\\
III &12 &$\ge0.999997944684$ &13 &$\ge0.999999999998$\\
V   &15 &$\ge0.927725586161$ &16 &$\ge0.999999964716$\\
\bottomrule
\end{tabular}
\end{center}
Here success means that the outer core is smooth and the solver's internal choices are good, conditional on $\Theta_F\not\equiv0$; it is at least
$(1-\beta_{V,s})(1-\alpha_{V,s})$.  At the safety fields, the base-two logarithms of $(\alpha_{V,s},\beta_{V,s})$ are
\[
(-27.994,-124.681),\qquad
(-39.858,-183.459),\qquad
(-24.756,-219.513).
\]
Thus the minimum fields are adequate for expected-time verified restart, while the safety fields give a clean one-invocation theorem.  The dependence on $\varepsilon$ in \eqref{eq:solver-extension-degree} is essential: no fixed list of extension degrees works for every arbitrarily small fixed $\varepsilon>0$.

\subsection{Published near-quadratic Kronecker solver}\label{subsec:solver}

\begin{proposition}[Fast smooth-core specialization]\label{prop:finite-field-kronecker}
Let $g_1,\ldots,g_V$ be homogeneous quadrics over a finite field $\F_{p^k}$.  Assume that they form a regular sequence and every projective prefix is geometrically radical.  Put $D_i=2^i$ and $D=D_V$.  Fix a rational $\varepsilon>0$, set $\eta=\varepsilon/3$, and assume
\[
 p^k>B_V(\eta).
\]
One finite invocation of Theorem~5.4 of van der Hoeven and Lecerf uses
\begin{equation}\label{eq:fast-kronecker}
 \widetilde O\!\left(D\,2^{(1+\varepsilon)(V-1)}k\log_2p\right)
\end{equation}
bit operations and fails because of its random coordinate, intersection, or evaluation choices with probability at most
\begin{equation}\label{eq:vdhl-internal-failure}
 \frac{5D^2+26D^3}{p^k}.
\end{equation}
Outside those bad choices it returns a complete Kronecker parametrization.  The checks in \cref{prop:resolution-checks}, including the requirement that the separator have degree $D$, certify completeness for this task.  Repeating checked finite invocations on an input satisfying the hypotheses gives the Las Vegas expected-time form of Corollary~5.5.

For $K=\F_{127^3}$ and any $s\ge s_\varepsilon^0$, the explicit finite-invocation cost over $L_s$ is obtained from \eqref{eq:fast-kronecker} with $k=3s$ and $p=127$.  Since $s=O(V)=O(\log D)$ for the choices above, this field-degree factor is absorbed by soft-$O$ notation.  Including field construction, casting, public slicing, output recombination, solving, and verification, the full expected projective route has scale
\begin{equation}\label{eq:fast-solver-scale}
 T_{\varepsilon,V}
 :=\widetilde O\!\left(
 2^V2^{(1+\varepsilon)(V-1)}3\log_2 127
 \right).
\end{equation}
\end{proposition}

\begin{proof}
The first statement is Theorem~5.4 specialized to the degree sequence $(2,\ldots,2)$, with internal parameter $\eta=\varepsilon/3$.  Its complete field condition is \eqref{eq:vdhl-field-bound}; its complexity exponent $1+3\eta$ becomes $1+\varepsilon$.  The three bad-choice bounds in the proof sum to \eqref{eq:vdhl-internal-failure}.

A smooth zero-dimensional projective core satisfies the required regularity and prefix-radicality assumptions by \cref{lem:full-implies-prefix}.  Corollary~5.5 of the source constructs a suitable extension, casts the input, verifies the returned parametrization, and obtains the same expected soft-$O$ scale.  Sampling $M$ and $R$, restricting the dense quadrics, and forming the output combinations costs only polynomially many operations in the working field and is lower order than the degree-dependent term.
\end{proof}

\begin{corollary}[One checked projective invocation]\label{cor:fast-single-invocation}
Fix $\varepsilon>0$, choose any $s\ge s_\varepsilon^0$, and fix a key with $\Theta_F\not\equiv0$.  Run one outer draw over $L_s$, one finite invocation from \cref{prop:finite-field-kronecker}, and all checks from \cref{prop:resolution-checks}; return \textsf{fail} if any check rejects.  The invocation has soft-$O$ cost $T_{\varepsilon,V}$ and succeeds with probability at least
\begin{equation}\label{eq:fast-single-invocation-success}
 (1-\beta_{V,s})(1-\alpha_{V,s}).
\end{equation}
For $\varepsilon=0.3$, the minimum-field and safety-field probabilities are those in the preceding table.
\end{corollary}

\begin{proof}
The outer discriminant fails with probability at most $\beta_{V,s}$.  Conditional on a smooth core, \cref{prop:finite-field-kronecker} fails internally with probability at most $\alpha_{V,s}$.  The two groups of coins are fresh.  Exact checking changes malformed or incomplete output into \textsf{fail}, never into a false accepted result.
\end{proof}

\subsection{A classical arithmetic cross-check}\label{subsec:slp-solver}

Let
\begin{equation}\label{eq:slp-length}
 M_2:=\binom{V+2}{2},
 \qquad
 L_{\rm quad}:=M_2+V(2M_2-1).
\end{equation}
The first term computes every degree-two monomial in the $V+1$ homogeneous variables once.  Each of the $V$ dense quadrics then uses $M_2$ scalar multiplications and $M_2-1$ additions.  Thus $L_{\rm quad}=\Theta(V^3)$ is an explicit division-free circuit for the full square core.

Giusti--Lecerf--Salvy give the older arithmetic bound
\begin{equation}\label{eq:classical-slp-scale}
 L_{\rm quad}\,\widetilde O(D^2)
\end{equation}
field operations under the same regularity conditions~\cite{GiustiLecerfSalvy2001,vdHL}.  We retain \eqref{eq:classical-slp-scale} only as an epsilon-free growth-rate cross-check.  We do \emph{not} transfer the explicit finite-field failure constants of Theorem~5.4 to this older bound: the literature cited here does not state that combined theorem.  Accordingly, the SLP line below has no proved one-invocation failure probability and is not used for the universal-forgery corollary.

\Needspace{16\baselineskip}
\begin{proposition}[Task-specific resolution checks]\label{prop:resolution-checks}
Let $f=(f_1,\ldots,f_V)$ be affine quadrics over a finite field $L$, and put $D=2^V$.  Suppose a proposed geometric resolution consists of a monic polynomial $m(T)$, coordinate polynomials $v_1(T),\ldots,v_V(T)$, and a linear form $u$.  Reject it unless every $v_j$ has degree below $\deg m$ and
\begin{enumerate}
\item $1\le\deg m\le D$ and $m$ is squarefree;
\item $u(v_1(T),\ldots,v_V(T))\equiv T\pmod m$;
\item $f_i(v_1(T),\ldots,v_V(T))\equiv0\pmod m$ for every $i$;
\item $\gcd\!\left(m(T),\det J_f(v_1(T),\ldots,v_V(T))\right)=1$.
\end{enumerate}
If the solver returns standard numerators $w_j$ with denominator $m'$, first verify $\gcd(m,m')=1$ and set $v_j=w_j(m')^{-1}\bmod m$.  Passing these checks certifies that the represented points are distinct simple solutions of the square core.  When this proposition is used with a smooth zero-dimensional complete intersection of $V$ quadrics, requiring in addition $\deg m=D$ certifies completeness by B\'ezout.  The checks cost $\widetilde O(\operatorname{poly}(V)D\log|L|)$ Boolean gates.
\end{proposition}

\begin{proof}
Squarefreeness gives distinct roots $\tau$ of $m$ over $\overline L$.  The primitive-form identity makes the represented points $P_\tau=(v_1(\tau),\ldots,v_V(\tau))$ distinct.  The equation checks put them in $V(f)$, and the Jacobian gcd makes each one simple.  Evaluate the quadrics and Jacobian in $L[T]/(m)$; compute the determinant there with a division-free algorithm because this quotient need not be a field.  Polynomial arithmetic, gcds, and the $V\times V$ determinant have the stated cost.
\end{proof}

\paragraph{Why the extra checks remain useful.}
The published Las Vegas solver already verifies its complete parametrization.  The task-specific checks make the cryptanalytic interface self-contained and protect against malformed implementations, truncated outputs, and conversion errors.  On every branch, final acceptance additionally requires all unrecombined equations, exact public-map factorization, and an invertible oil system.  No malformed solver transcript can create a false accepted key.

\subsection{Root selection and oil-space recovery}\label{subsec:root-recovery}

\begin{lemma}[Recovery from the projective oil point]\label{lem:descent}
Let $L/K$ be a finite extension.  Suppose $\Theta_F(M,R)\ne0$, and let $[t^\star]\in\PP^V(L)$ be the unique point for which $x^\star=Mt^\star$ lies in $\PP(\cO)_L$.  Then
\[
\rank(P_1x^\star\ \cdots\ P_mx^\star)=V,
\qquad
\ker(P_1x^\star\ \cdots\ P_mx^\star)^{\mathsf T}=\cO_L.
\]
Row reduction recovers the public graph of the oil space.  For the planted point that graph is defined over $K$ and, after the specification's inverse pull-back, yields the planted expanded UOV decomposition.
\end{lemma}

\begin{proof}
The full-rank matrix $M$ defines a projective $V$-plane over $L$.  Its vector-space intersection with the $O$-dimensional oil space is nonzero by dimension count.  If that intersection had dimension at least two, all restricted quadrics would vanish on a positive-dimensional projective linear space, contradicting smooth zero-dimensionality.  Thus it is one $L$-rational point.

Smoothness of the square core gives derivative rank at least $V$ there.  Total isotropy puts the $O$-dimensional oil space inside the ambient common derivative kernel, so the rank is at most $(V+O)-O=V$.  Equality holds, and \cref{lem:kernel} identifies the kernel with $\cO_L$.  The remaining statements follow from $K$-descent and \cref{prop:verified-key}.
\end{proof}

\begin{lemma}[Separator descent]\label{lem:separator-descent}
Let a reduced zero-dimensional scheme over $L=\F_{\mathfrak q}$ have a Kronecker representation whose primitive linear form $u$ separates its geometric points.  A geometric point $P$ is $L$-rational if and only if $u(P)\in L$.
\end{lemma}

\begin{proof}
The forward implication is immediate.  Conversely, if $u(P)\in L$, then
\[
u(P^{(\mathfrak q)})=u(P)^{\mathfrak q}=u(P).
\]
Because $u$ is defined over $L$ and separates the geometric points, Frobenius must fix $P$.
\end{proof}

\Needspace{15\baselineskip}
\begin{proposition}[Filter all original equations before factoring]\label{prop:separator-prefilter}
Let a consistency-checked resolution over $L_s$ have squarefree separator polynomial $m_u(T)$ and coordinate polynomials $v(T)$.  Let $\widetilde q_1,\ldots,\widetilde q_m$ be the dehomogenized original restricted quadrics in the same affine chart, and define
\begin{equation}\label{eq:all-equation-factor}
a(T):=\gcd\bigl(m_u(T),\widetilde q_1(v(T)),\ldots,\widetilde q_m(v(T))\bigr).
\end{equation}
Then the roots of $a$ are exactly the separator values of represented points satisfying all $m$ original equations.  Moreover,
\begin{equation}\label{eq:rational-factor}
a_{L_s}(T):=\gcd\bigl(a(T),T^{q_s}-T\bigr)
\end{equation}
has exactly the separator values of those common points that are $L_s$-rational.  Computing $a$ and $a_{L_s}$ costs $\softO\!\left(D\,\operatorname{poly}(V,m,\log|L_s|)\right)$ Boolean gates.  Only $a_{L_s}$, whose degree may be much smaller than $D$, must be factored.
\end{proposition}

\begin{proof}
Work over an algebraic closure of $L_s$.  Because $m_u$ is squarefree, it is the product of the distinct factors $T-\tau$ indexed by the represented points.  For each root $\tau$,
\[
\widetilde q_i(v(T))\bmod m_u\;\big|_{T=\tau}=\widetilde q_i(v(\tau)).
\]
Thus $T-\tau$ divides \eqref{eq:all-equation-factor} exactly when every original restricted equation vanishes at the represented point.  The second claim follows from \cref{lem:separator-descent} and the identity $z^{q_s}=z$ for $z\in L_s$.  Dense quadratic evaluation in $L_s[T]/(m_u)$, modular exponentiation, and polynomial gcd computation at degree at most $D$ give the stated cost.  Since $\log|L_s|=O(V)$ for the selected fields, it is lower order than the degree-squared solver bound.
\end{proof}

\paragraph{Check the resolution, prefilter, and extract rational roots.}
Convert the solver output to canonical coordinate polynomials and run the consistency checks of \cref{prop:resolution-checks}.  On failure, discard the invocation.  On a successful Las Vegas solver call, the returned separator polynomial $m_u(T)\in L_s[T]$ represents the full affine zero scheme and contains the planted oil point.  The additional checks reject malformed or incomplete implementation outputs.

Apply \cref{prop:separator-prefilter}: evaluate all $m$ original restricted equations in $L_s[T]/(m_u)$, compute $a(T)$, and then compute $a_{L_s}(T)$.  Factor only $a_{L_s}$ over $L_s$.  A Las Vegas finite-field factorization algorithm terminates with probability one and has expected subquadratic dependence on the input degree, uniformly for inputs of degree at most $D$; its expected contribution is therefore below the solver's near-quadratic-in-$D$ scale~\cite{KaltofenShoup}.  Reconstruct the corresponding affine points, undo the coordinate changes, and recheck all $m$ original equations as a final exact verification.

\paragraph{Certify the decomposition.}
For each survivor $x$, compute the common derivative kernel.  Require an $O$-dimensional graph defined over $K$, total isotropy, and every identity \eqref{eq:graphid}.  Apply the specification's inverse pull-back and verify the base-field factorization against the original public map.  An output is called an \emph{exact candidate} only after all these checks pass.

\paragraph{Certify signing or restart.}
Deduplicate exact candidates by their $K$-rational graph.  For each remaining candidate, first test the $V$ basis assignments from \cref{prop:signing-search}.  If they all fail, scan the remaining nonzero combinations on the designated pair lines.  Return the first candidate with an invertible oil matrix, together with that matrix's inverse.  If no candidate passes, discard the trial and sample a fresh slice.  On a correct solver branch the candidate list contains the planted decomposition, so outside the event in \eqref{eq:eta-sign} this scan succeeds deterministically.  At most
\[
 D\bigl(V+\lfloor V/2\rfloor(q-1)\bigr)
\]
oil matrices are tested per trial, below the solver's near-quadratic-in-$D$ scale.  In the public-random model, the basis-first part succeeds after fewer than $1.008$ tests on average.

Let $\mathsf A_{\varepsilon,V,s}$ denote the event that the full trial accepts a publicly verified signing decomposition and inverse oil matrix.  On a key outside $\cE_{\rm sign}$, a draw with $\Theta_F(M,R)\ne0$ yields the planted exact candidate whenever the finite solver invocation succeeds, after which the deterministic signing scan succeeds.  The outer and internal coins are fresh, so
\begin{equation}\label{eq:acceptance-lower-bound}
\Pr[\mathsf A_{\varepsilon,V,s}]
\ge(1-\beta_{V,s})(1-\alpha_{V,s}).
\end{equation}
This one-sided implication is the only interface needed between the solver guarantee and the verified attack.  On irregular inputs or bad internal choices, a failed check returns \textsf{fail}; it never creates a false accepted key.

\begin{lemma}[Verified-restart accounting]\label{lem:restart-accounting}
Suppose each fresh trial, conditional on any history of earlier failures, accepts with probability at least $p>0$.  Then the number $N$ of trials satisfies
\[
\Pr[N>t]\le(1-p)^t,
\qquad
\mathbb E[N]\le p^{-1}.
\]
If the current trial uses fresh coins and has expected work at most $T$, uniformly over prior histories, then the expected total work is at most $T/p$.
\end{lemma}

\begin{proof}
The tail bound follows inductively from the conditional success lower bound, and summing the tail probabilities gives $\mathbb E[N]\le1/p$.  For the work bound, the event that trial $j$ is reached depends only on earlier trials.  Fresh current-trial coins therefore give
\[
\mathbb E\!\left[\mathbf 1_{\{N\ge j\}}W_j\right]
\le T\Pr[N\ge j].
\]
Summing over $j$ and applying the trial bound proves the claim.
\end{proof}

\Needspace{15\baselineskip}
\subsection{Main recovery theorem}\label{subsec:main-theorem}
Use $T_{\varepsilon,V}$ from \eqref{eq:fast-solver-scale}.  For $0<\epsilon<1/4$ and an extension degree $s$ satisfying \eqref{eq:local-field-condition}, write
\[
 T_{\rm loc}(\epsilon,s):=
 \widetilde O\!\left(
 \bigl(V^4+V(L_{\rm loc}+V^2+V^4)\bigr)V4^V\log_2(Q^s)
 \right).
\]

\begin{theorem}[Key-by-key recovery from one graph direction]\label{thm:graph-attack}
Let an expanded \QRUOV{} key over $K$ have a valid planted decomposition.  Suppose that at least one maximal minor of $\mathcal B(c)$ is a nonzero polynomial in $c$.  Fix $0<\epsilon<1/4$ and choose $s$ so that \eqref{eq:local-field-condition} holds.  Put
\[
 p_{\rm dir}:=\left(1-\frac VQ\right)
 \left(1-\frac{V(m-V)}Q\right)(1-2\epsilon).
\]
Then the following verified attacks are correct for this fixed key.
\begin{enumerate}
\item A randomized graph-direction trial finds an exactly verified expanded UOV decomposition with probability at least $p_{\rm dir}$.  Repeating fresh trials terminates almost surely, uses at most $p_{\rm dir}^{-1}$ trials on average, and has expected work
\begin{equation}\label{eq:local-expected-decomposition}
 O\!\left(T_{\rm loc}(\epsilon,s)/p_{\rm dir}\right).
\end{equation}
Every accepted output is correct.

\item If the determinant polynomial $\Delta$ of the planted decomposition is nonzero, testing one fresh uniform scalar-block vinegar assignment in each trial returns a signing certificate with expected trial count at most
\begin{equation}\label{eq:local-functional-trials}
 \bigl(p_{\rm dir}(1-m/q)\bigr)^{-1}.
\end{equation}

\item If one planted basis matrix is invertible, the deterministic basis-first signing scan succeeds on every successful recovery trial.  Under the designated-pair condition of \cref{prop:signing-search}, the full basis-and-line scan does so after at most $V+\lfloor V/2\rfloor(q-1)$ oil-matrix tests per candidate.  In either case the expected trial and work bounds are those of item~1.

\item If a coordinate direction $e_j$ satisfies $\rank\mathcal B(e_j)=V$, all outer attack choices can be made deterministic: scan the $O\,[V(m-V)+1]$ direction/condenser pairs of \cref{cor:deterministic-outer-list}, and repeat the finitely many probabilistic solver calls in rounds.  This still terminates almost surely and adds only a polynomial factor.  The shorter lists of \cref{cor:compact-direction-list} give category-scale failure bounds in the public-random model.
\end{enumerate}
The theorem uses only a local nonsingularity condition at the planted root.  It does not require the projective discriminant, the off-oil incidence argument, a key-generation distribution, or uniqueness of the recovered decomposition.
\end{theorem}

\begin{proof}
\proofstep{A good trial contains the planted root.}
By \cref{prop:graph-direction-density}, a random $c$ makes $\mathcal B(c)$ full rank with probability at least $1-V/Q$.  Conditional on that event, a random $t$ makes $\mathsf C(t)\mathcal B(c)$ invertible with probability at least $1-V(m-V)/Q$.  The planted root then lies in the localized smooth set by \cref{cor:local-hypotheses}.  By \cref{cor:gimenez-local}, the solver returns a representation containing it with probability at least $1-2\epsilon$.

\proofstep{One planted root gives the whole hidden graph.}
For every returned root $x$, solve the public systems \eqref{eq:direction-completion-system}.  At $x=Gc$, \cref{prop:graph-direction-completion} returns every column of $G$.  The attack then checks all original direction equations, rank conditions, graph identities, inverse pull-back identities, and public-map factorizations.  Any root that does not yield a valid graph is discarded.

\proofstep{Restart and signing.}
A fresh trial therefore reaches the planted expanded decomposition with probability at least $p_{\rm dir}$.  Apply \cref{lem:restart-accounting}.  The randomized and deterministic signing claims follow from \cref{prop:signing-search}.  Filtering and completing at most $D$ roots costs $\widetilde O(D\operatorname{poly}(V,O,m,\log|L|))$, below \eqref{eq:local-solver-cost}.  Exact verification prevents false acceptance even when the internal Monte Carlo solver fails.

For item~4, \cref{cor:deterministic-outer-list} guarantees that a regular coordinate direction has a condenser value with nonsingular planted Jacobian.  The deterministic scan reaches it.  Repeating all finitely many solver calls in rounds gives a fixed positive success probability per round and hence almost-sure termination.
\end{proof}

\begin{corollary}[Public-random density and compact deterministic certificates]\label{cor:graph-ideal-recovery}
In the public-random expansion model, the probability that all $O$ coordinate directions are rank deficient is exactly $\tau_{m,V}(Q)^O$, whose base-two logarithms are
\[
 -1132.166907,\qquad-1635.352198,\qquad-2935.247544.
\]
Hence, except with probability at most
\begin{equation}\label{eq:eta-local-functional}
 \eta_{\rm local}:=\tau_{m,V}(Q)^O+\eta_{\rm sign},
\end{equation}
the full deterministic direction list and deterministic signing fallback recover a universal signing certificate.

A much shorter fixed attack already has category-scale coverage.  Use the first $3,4,4$ coordinate directions and their condenser lists, together with the first $22,31,40$ basis signing assignments, at Levels I, III, and V.  These attacks solve only $315,612,1228$ square systems and test at most $22,31,40$ oil matrices per exact candidate.  Their failure probabilities are at most
\[
 \tau_{m,V}(Q)^{r_{\rm dir}}+\sigma_m^{b_{\rm sign}},
\]
with base-two logarithms below
\begin{equation}\label{eq:compact-combined-logs}
 -153.502,\qquad-216.298,\qquad-279.095.
\end{equation}
No independence between the recovery and signing failures is needed for this union bound.
\end{corollary}

\begin{proof}
The all-direction statement is \eqref{eq:all-coordinate-directions-bad}, followed by \cref{prop:signing-search} and the union bound.  For the compact statement, combine \cref{cor:compact-direction-list} with the fixed-prefix signing bound \eqref{eq:compact-sign-logs}.  On the complement, the deterministic outer list contains a square system with the planted point nonsingular, and the fixed signing list contains an invertible planted oil matrix.  Round-robin solver repetition and exact verification then give the claimed universal certificate.
\end{proof}

\begin{theorem}[Fast projective key-by-key recovery]\label{thm:attack}
Let an expanded \QRUOV{} key over $K$ have a valid planted expanded decomposition, and assume that the public joint slice-and-output discriminant $\Theta_F(M,R)$ is not the zero polynomial.  Fix a rational $\varepsilon>0$, choose any $s\ge s_\varepsilon^0$, and put
\[
 p_{\varepsilon,V,s}:=(1-\beta_{V,s})(1-\alpha_{V,s}).
\]
Then:
\begin{enumerate}
\item A decomposition-only verified-restart algorithm terminates with probability one and returns an exactly verified expanded UOV decomposition.  Every accepted output is correct, the expected number of outer trials is at most $p_{\varepsilon,V,s}^{-1}$, and the expected work is
\begin{equation}\label{eq:expected-decomposition}
 O\!\left(\frac{T_{\varepsilon,V}}{p_{\varepsilon,V,s}}\right).
\end{equation}
For $\varepsilon=0.3$, the minimum fields are $s=8,12,15$.  Their expected-trial bounds are below
\[
1.00773,\qquad1.000003,\qquad1.07791.
\]
The safety fields $s=9,13,16$ give the near-one single-invocation probabilities in \cref{cor:fast-single-invocation}.

\item Suppose the planted determinant polynomial $\Delta$ from \cref{prop:signing-search} is nonzero.  If each outer trial tests one fresh uniform scalar-block vinegar assignment for every exact candidate, verified restart returns an equivalent signing key and inverse oil matrix.  The expected number of outer trials is at most
\[
 \frac{1}{p_{\varepsilon,V,s}(1-m/q)},
\]
and the expected work is
\begin{equation}\label{eq:expected-attack-randomized}
 O\!\left(\frac{T_{\varepsilon,V}}{p_{\varepsilon,V,s}(1-m/q)}\right).
\end{equation}
At most $D$ oil matrices are tested per trial.

\item If one planted basis matrix $M_j$ is invertible, the basis-first deterministic scan returns an equivalent signing key on every successful core invocation and tests at most $DV$ oil matrices per outer trial.  Under the designated-pair condition of \cref{prop:signing-search}, the full deterministic fallback tests at most
\[
 D\bigl(V+\lfloor V/2\rfloor(q-1)\bigr)
\]
oil matrices per trial.  In either case the expected trial count is at most $p_{\varepsilon,V,s}^{-1}$ and the expected work remains $O(T_{\varepsilon,V}/p_{\varepsilon,V,s})$.
\end{enumerate}
These are statements about each fixed expanded key satisfying the displayed conditions.  They do not depend on how the key was generated.  Recovering the original seed or serialized private key is unnecessary.
\end{theorem}

\begin{proof}
\proofstep{A regular draw yields a complete core resolution.}
The discriminant hypothesis and \eqref{eq:fast-beta} imply that a fresh $(M,R)$ gives a smooth full core with probability at least $1-\beta_{V,s}$.  \Cref{lem:full-implies-prefix} supplies the solver hypotheses.  Conditional on a smooth core, the finite invocation in \cref{prop:finite-field-kronecker} returns and verifies the complete geometric resolution with probability at least $1-\alpha_{V,s}$.  Every invocation has a finite operation schedule and returns \textsf{fail} if a denominator, degree, or verification check rejects.  An irregular outer draw therefore cannot trap the attack inside an unjustified restart loop.

\proofstep{The planted decomposition appears among the exact candidates.}
On a regular draw, the quotient-algebra prefilter of \cref{prop:separator-prefilter}, rational-root extraction, and \cref{lem:descent} retain the planted oil point and recover the planted expanded decomposition.  Exact all-equation checks, inverse pull-back, and public-map factorization reject every invalid candidate.

\proofstep{Signing certification.}
If $\Delta\not\equiv0$, \cref{prop:signing-search} gives success probability at least $1-m/q$ for a fresh uniform assignment on the planted candidate.  If a basis matrix or designated pair witness exists, the corresponding deterministic scan succeeds on every regular draw.

\proofstep{Restart and auxiliary work.}
Apply \cref{lem:restart-accounting}.  Finite-field factorization has subquadratic expected dependence on degree, and all filtering, derivative-kernel, descent, pull-back, verification, and signing work costs $\widetilde O(D\operatorname{poly}(V,O,m,q,\log|L_s|))$.  For every fixed $\varepsilon>0$ this is lower order than $T_{\varepsilon,V}$.  Exact verification prevents false acceptance on every branch.
\end{proof}

\begin{corollary}[Public-random density for the fast projective route]\label{cor:ideal-recovery}
In the public-random expansion model, the probability that either $\Theta_F\equiv0$ or the planted determinant polynomial $\Delta$ is identically zero is below
\[
2^{-364.563},\qquad2^{-457.919},\qquad2^{-1038.066}
\]
at Levels I, III, and V.  Every other key satisfies the randomized equivalent-key conclusion of \cref{thm:attack}.  In fact, outside the same bounded union of bad events, one designated pair span contains an invertible matrix, so the deterministic refinement applies.
\end{corollary}

\begin{proof}
The event $\Theta_F\equiv0$ is bounded by \cref{thm:key-density}.  If $\Delta$ is the zero polynomial, every designated pair span is singular, so this event is contained in $\cE_{\rm sign}$ and is bounded by \cref{prop:signing-search}.  The union has probability at most
\[
\eta_{\off}+\eta_{\oil}+\eta_{\rm span}+\eta_{\rm sign},
\]
which gives the displayed numerical bounds.  On the complement of this particular union, the designated-pair condition itself holds, so item~3 of \cref{thm:attack} applies.
\end{proof}

\begin{corollary}[One-invocation key-only universal forger]\label{cor:single-invocation-forger}
In the public-random expansion model, set $\varepsilon=0.3$, use the safety field $s=s_{0.3}^+$, run the checked projective invocation of \cref{cor:fast-single-invocation}, and then run the basis-first scan and projective-line fallback of \cref{prop:signing-search}.  The adversary uses no signing oracle, never outputs an invalid signature, and can forge every chosen message whenever it succeeds.  Its success probability over the key and attack coins is at least
\begin{equation}\label{eq:one-pass-forger-success}
 1-\delta_{\rm key}-\alpha_{V,s_{0.3}^+}-\beta_{V,s_{0.3}^+},
\end{equation}
where $\delta_{\rm key}$ is the key-level exceptional probability in \cref{cor:ideal-recovery}.  In particular, the success probabilities at Levels I, III, and V are respectively at least
\[
 0.999999996259,\qquad0.999999999998,\qquad0.999999964716.
\]
\end{corollary}

\begin{proof}
Outside the key-level exceptional set, the projective discriminant and a deterministic pair-line witness both exist.  A successful checked solver pass returns the complete core, so the planted point survives exact filtering and yields the planted expanded decomposition.  The deterministic signing scan then returns a certificate that inverts every public target.  Apply \cref{cor:universal-forgery} and union-bound the key, outer-discriminant, and internal-solver failures.  Exact verification ensures that failure affects only whether a forgery is returned, never whether a returned forgery is valid.
\end{proof}

\section{Asymptotic cost and parameter implications}\label{sec:security}

\subsection{What the cost numbers mean}

The QR-UOV specification reports classical security estimates as Boolean-gate exponents~\cite[Section~6.3 and Table~14]{QRUOVSpec}.  Our main solver result is an asymptotic bit-complexity theorem, so it is reasonable to compare exponents at the same high level.  It is not reasonable to interpret the result as a measured gate count.  For a fixed rational $\varepsilon>0$, the displayed leading term gives the normalized soft-$O$ exponent
\begin{equation}\label{eq:fast-indicator}
 C_\varepsilon(V)
 :=V+(1+\varepsilon)(V-1)+\log_2(3\log_2 127).
\end{equation}
The working-field construction and casting are already included in the theorem.  For the primary choice $\varepsilon=0.3$, using the minimum admissible fields changes the base-two expected-work exponent by at most $0.109$ bits; using the one-degree safety fields changes it by at most $5.1\times10^{-8}$ bits.

For a convention independent of the solver's internal circuit presentation, we also report
\begin{equation}\label{eq:uniform-indicator}
 U_\varepsilon(V):=C_\varepsilon(V)+\log_2 C_\varepsilon(V),
\end{equation}
which exposes one unit-coefficient $T\log T$ simulation of a $T$-step bit algorithm~\cite{PippengerFischer}.  The specification does not apply both its field-operation conversion and a second generic simulation, so the two columns are sensitivity views rather than additive charges.

Neither \eqref{eq:fast-indicator} nor \eqref{eq:uniform-indicator} is a complete circuit estimate.  Both suppress constants and polylogarithmic factors, and the hidden constant may depend badly on $\varepsilon$.  The tables therefore answer a narrower question: whether the published asymptotic algorithm has enough exponent margin to justify a closer concrete analysis.  They do not, by themselves, certify that a NIST category target has been crossed.

\subsection{Sensitivity to the solver exponent}

\begin{table}[t]
\centering
\small
\caption{Fast smooth-core line.  Each cell is direct $C_\varepsilon(V)$ / uniform $U_\varepsilon(V)$.}
\label{tab:eps-sensitivity}
\setlength{\tabcolsep}{4pt}
\begin{tabular}{@{}c r r r r r@{}}
\toprule
Level & $\varepsilon=0.1$ & $\varepsilon=0.2$ & $\varepsilon=0.3$ & $\varepsilon=0.5$ & Target \\
\midrule
I   & $112.49/119.30$ & $117.59/124.47$ & $122.69/129.63$ & $132.89/139.94$ & $143$ \\
III & $162.89/170.24$ & $170.39/177.80$ & $177.89/185.36$ & $192.89/200.48$ & $207$ \\
V   & $217.49/225.25$ & $227.59/235.42$ & $237.69/245.58$ & $257.89/265.90$ & $272$ \\
\bottomrule
\end{tabular}
\end{table}

The paper uses $\varepsilon=0.3$ as its primary fixed choice.  It is not a formal limit, yet its uniform row lies $13.371,21.635,26.417$ bits below target.  The much more conservative $\varepsilon=0.5$ row remains below target by $3.056,6.518,6.099$ bits.

\begin{proposition}[Robustness envelope]\label{prop:eps-robust}
For the three recommended dimensions, the normalized indicator $U_\varepsilon(V)$ is below the nominal target whenever
\[
 \varepsilon<0.559284,\qquad0.586277,\qquad0.560058
\]
at Levels I, III, and V, respectively.
\end{proposition}

\begin{proof}
The function $U_\varepsilon(V)$ is continuous and strictly increasing in $\varepsilon$.  Solving $U_\varepsilon(V)$ equal to the target gives the three thresholds.  The accompanying audit recomputes every value from \eqref{eq:fast-indicator}--\eqref{eq:uniform-indicator}.
\end{proof}

\begin{remark}[Margins are not certified overhead reserves]\label{rem:no-overhead-reserve}
For fixed $\varepsilon$, subtracting $U_\varepsilon(V)$ from the target is a useful comparison of normalized exponents.  It is not a theorem that the solver can absorb an arbitrary multiplicative factor of that size: the constant hidden by $\widetilde O(\cdot)$ is unspecified and may itself depend on $\varepsilon$.  We therefore report bit margins, but do not call their exponentials ``whole-attack budgets.''
\end{remark}

\subsection{Epsilon-free arithmetic cross-check}\label{subsec:slp-cost}

The older straight-line-program theorem gives the arithmetic scale
$L_{\rm quad}\widetilde O(D^2)$ from \eqref{eq:classical-slp-scale}.  To see what that scale means at the recommended dimensions, define the diagnostic operand-size parameter
\[
 s_{\rm diag}:=\min\{s\ge1:Q^s\ge100D^3\}
\]
and
\begin{align}
 A_0(V)&:=2V+\log_2L_{\rm quad}+\log_2(s_{\rm diag}\log_2Q),\label{eq:slp-direct-indicator}\\
 A_1(V)&:=A_0(V)+\log_2A_0(V).\label{eq:slp-stress-indicators}
\end{align}
The first row assigns unit coefficient to the displayed arithmetic term after charging for arithmetic in an $s_{\rm diag}$-degree extension.  The second adds one generic $T\log T$ sensitivity.  These are arithmetic-growth diagnostics only: \cref{subsec:slp-solver} explains why we do not attach the modern theorem's exact one-invocation failure constants to this older bound.

\begin{table}[t]
\centering
\small
\caption{Epsilon-free classical SLP arithmetic cross-check.  Margins are target minus indicator, in bits.}
\label{tab:slp-costs}
\setlength{\tabcolsep}{5pt}
\begin{tabular}{@{}c r r r r r@{}}
\toprule
Level & $s_{\rm diag}$ & $A_0$ & Margin & $A_1$ & Margin \\
\midrule
I   & 8  & $128.587$ & $14.413$ & $135.593$ & $7.407$ \\
III & 12 & $178.784$ & $28.216$ & $186.266$ & $20.734$ \\
V   & 15 & $232.363$ & $39.637$ & $240.223$ & $31.777$ \\
\bottomrule
\end{tabular}
\end{table}

The diagnostic agrees with the modern theorem on the dominant growth: about $4^V$ times polynomial factors.  Its one-log sensitivity remains below all three targets, but the table is not a finite-gate attack estimate and has no standalone success-probability claim.  Its role is to show that the below-target trend is not created solely by choosing a favorable value of $\varepsilon$ in the newer theorem.

\begin{proposition}[Arithmetic growth check]\label{prop:slp-growth}
For fixed $Q$,
\[
 \log_2\!\left(L_{\rm quad}D^2s_{\rm diag}\log_2Q\right)
 =2V+4\log_2V+O(1).
\]
Consequently, within this arithmetic model, any parameter family targeting $\lambda$ classical bits and keeping $V\le(\tfrac12-\delta)\lambda$ for a fixed $\delta>0$ is eventually below target, even after multiplication by any fixed polynomial in $V$.
\end{proposition}

\begin{proof}
The circuit length in \eqref{eq:slp-length} is $\Theta(V^3)$, while $s_{\rm diag}=\Theta(V)$ and $D=2^V$.  Taking base-two logarithms gives $2V+3\log_2V+\log_2V+O(1)$.  A fixed polynomial multiplier contributes only $O(\log\lambda)$ bits, whereas the assumed gap to $\lambda$ grows linearly.
\end{proof}

\subsection{Separate explicit-prefactor solver lines}\label{subsec:independent-cost}

The modern fast line has a fixed solver exponent $\varepsilon$ and a hidden soft-$O$ constant; the classical SLP line removes $\varepsilon$ but remains in the same Kronecker lineage.  To separate the cryptanalytic reduction from both of those interfaces, we also specialize \cref{prop:gimenez-theorem} directly.  For each level, the audit chooses $\epsilon=2^{-b}$ and the smallest extension degree $s$ satisfying the theorem's field-size condition.  The ``direct'' column takes the displayed polynomial prefactor in \eqref{eq:gimenez-cost} with unit coefficient and includes the expected retry factor $(1-2\epsilon)^{-1}$.  The ``stress'' column adds one generic $T\log T$ simulation factor.  Both remain normalized soft-$O$ indicators, not measured circuits.

\begin{table}[t]
\centering
\small
\caption{Independent Gim\'enez-theorem specializations.  Each cost cell is direct / generic $T\log T$ stress, in base-two bits.}
\label{tab:gimenez-costs}
\setlength{\tabcolsep}{4pt}
\begin{tabular}{@{}c c c r r r@{}}
\toprule
Level & Route & $(\epsilon,s)$ & Direct & Stress & Target\\
\midrule
I   & smooth projective & $(2^{-13},9)$ & $141.119$ & $148.260$ & $143$\\
    & graph local       & $(2^{-8},9)$  & $146.803$ & $154.000$ & \\
III & smooth projective & $(2^{-23},13)$& $192.369$ & $199.956$ & $207$\\
    & graph local       & $(2^{-18},13)$& $198.597$ & $206.231$ & \\
V   & smooth projective & $(2^{-7},16)$ & $246.803$ & $254.751$ & $272$\\
    & graph local       & $(2^{-22},17)$& $253.526$ & $261.512$ & \\
\bottomrule
\end{tabular}
\end{table}

The smooth-projective direct row is below every target, although its generic stress exceeds Level I.  The graph-local route deliberately pays for a degree-$V$ localization polynomial and a division-free determinant circuit; its direct and stress rows remain below the Level-III and Level-V targets.  Thus the concern at Levels III and V is visible through both the projective and graph-local reductions and through two published solver theorems.  No analogous concrete claim is made for Level I.

For transparency, the local table uses the normalized determinant allowance $L_{\rm loc}=L_{\rm core}+V^4$ suggested by the $O(V^4)$ Samuelson--Berkowitz bound.  Replacing that allowance by $L_{\rm core}+4V^4$ changes the local direct/stress rows to
\[
 148.107/155.318,\qquad199.908/207.551,\qquad254.840/262.833.
\]
This harsher sensitivity still leaves the Level-V line well below target and the Level-III direct line below target; it is not used to make a Level-I claim.  The purpose of the local route is to remove the global-geometric single point of failure, not to manufacture the smallest numerical row.

\subsection{Comparison with the Round-2 rows}

The Round-2 specification reports classical attack rows $150,211,277$ bits and category targets $143,207,272$ for Levels I, III, and V~\cite[Table~14]{QRUOVSpec}.  For the primary choice $\varepsilon=0.3$, the modern uniform indicators are $20.37,25.64,31.42$ bits below the reported attack rows and $13.37,21.64,26.42$ bits below the targets.  The epsilon-free SLP one-log diagnostics are $14.41,24.73,36.78$ bits below the reported rows and $7.41,20.73,31.78$ bits below the targets.

The same pattern also appears through the independent Gim\'enez theorem.  Its smooth-projective direct row is below every target; its generic $T\log T$ stress row is below the Level-III and Level-V targets but above Level I.  The graph-local route, which needs only a nonsingular planted root rather than a global projective smoothness statement, is also below the Level-III and Level-V targets in the displayed direct calculation.  Thus Levels III and V show substantial margin through several theorem interfaces.  Level I remains more sensitive to hidden constants and representation choices.

These are asymptotic warning signs, not replacement concrete attack rows.  The calculations combine published solver bounds with the actual QR-UOV dimensions and field sizes, but do not measure the hidden constants, low-level arithmetic, or memory traffic of a full implementation.  Their value is that the structural reduction, the modern near-quadratic theorem, the older arithmetic growth check, and the independent locally closed-set theorem all point in the same direction.

\paragraph{Parameter implications.}
The results justify a concrete re-evaluation of all three parameter sets, with the strongest evidence at Levels III and V.  Such a review should implement or tightly upper-bound the dominant solver stages, include memory and low-level arithmetic costs, and then jointly reassess $V,m,O$, key size, signing failure, and the other known attack families.  This paper does not propose replacement parameter sets.

\paragraph{What remains unproved.}
The tables do not show that a full-parameter implementation uses fewer than $2^{143},2^{207},2^{272}$ physical gates.  The direct solver also requires astronomical space at the recommended dimensions, as \cref{subsec:memory} makes explicit.  The rigorous claims are the algebraic recovery reductions, their key-density statements in the stated models, the exact verification and failure accounting, and the cited asymptotic solver bounds.  A concrete category break requires an additional implementation-level cost analysis.

\subsection{Seeded expansion}

The key-by-key theorem itself needs no idealization.  To obtain a high-probability theorem for seeded key generation, only the conditional uniformity of the $m$ public rectangular streams and the product-uniform marginal of the $m$ symmetric streams must be compared with the public-random model; arbitrary graph--symmetric-block correlation is allowed.  \Cref{app:seeded-transfer} proves the needed comparison, with explicit statistical losses, when SHAKE is modeled as an ideal XOF or AES as an ideal cipher.  Adding those losses to the much smaller ideal-key exceptions gives the following base-two logarithms:
\[
\begin{array}{c|rrr}
&\text{Level I}&\text{Level III}&\text{Level V}\\ \hline
\text{SHAKE ideal XOF}&-121.245&-184.715&-248.286\\
\text{AES ideal cipher}&-99.863&-96.655&-94.112
\end{array}
\]
On every other seeded key in the corresponding ideal-primitive model, the same verified-restart attack terminates almost surely with the same key-by-key expected-work bound and acceptance with no false positives.  No factorization timeout or fixed attack budget is needed: the statistical comparison is applied only once, to the key-level exceptional set.

This does not yield a high-probability key-generation theorem for the fixed SHAKE/AES primitives.  It does not weaken the key-by-key theorem: every fixed-primitive key satisfying those key-by-key conditions is already covered.  A deterministic $\lambda$-bit expansion has support at most $2^\lambda$, whereas the independent-stream experiment has vastly larger support; \cref{prop:seed-support-barrier} makes the resulting near-maximal total-variation gap explicit.  Any standard-model density theorem must therefore analyze the particular exceptional event or assume computational pseudorandomness.  These observations delimit the fixed-primitive probability gap; they do not affect the key-by-key, ideal-XOF, or ideal-cipher theorems.

\subsection{Memory}\label{subsec:memory}

The cost comparisons above count Boolean operations, not storage.  Here that distinction is decisive.  A smooth square core has degree $D=2^V$, so a conventional univariate resolution contains degree-$D$ polynomials.  Even one dense vector of $D$ elements over the minimum solver field is enormous.

For the minimum $\varepsilon=0.3$ fields $s=8,12,15$, a random affine chart contains all geometric roots with probability at least $1-D/|L_s|$.  The logarithms of the chart-miss bounds are $-115.728,-175.593,-212.491$.  Storing one degree-$D$ vector with $s\log_2Q$ bits per field element requires at least
\[
0.0819\ \mathrm{EiB},\qquad
2.06\times10^6\ \mathrm{EiB},\qquad
1.73\times10^{14}\ \mathrm{EiB}
\]
at Levels I, III, and V.  A dense primitive-element representation stores roughly $V$ such coefficient arrays, requiring about
\[
4.26\ \mathrm{EiB},\qquad
1.57\times10^8\ \mathrm{EiB},\qquad
1.76\times10^{16}\ \mathrm{EiB},
\]
before temporary workspace.  The one-degree safety fields are larger still.

\Cref{prop:implicit-root,prop:resultant-gradient,rem:implicit-memory} show that the desired \emph{output} can be much smaller: one isolated root can be encoded by $V+1$ norm derivatives, one resultant-gradient block, or a one-dimensional quotient-algebra eigenspace.  With black-box quotient multiplication, the eigenspace route needs only one $O(D)$ state vector rather than all $V$ coordinate arrays.  This saves a factor of $V$ in post-solver storage, but it does not show how to construct the quotient algebra in low space.  The published direct solvers are therefore asymptotically informative but not implementable at the recommended parameters.

\section{Conclusion}\label{sec:conclusion}

One regular public direction is enough to recover the hidden oil space.  For a chosen $c\in K^O$, the public system
\[
 P_i(-x,c)=0\qquad(1\le i\le m)
\]
contains the planted point $x=Gc$.  When the derivative matrix at that point has rank $V$, solving for this one point turns recovery of the remaining columns of $G$ into public linear algebra.  Exact graph identities, pull-back checks, and public-map factorization make false acceptance impossible.

This local reduction is stronger and simpler than the original projective route.  It needs no global smoothness statement and no uniqueness assumption.  In the public-random model, the first $3,4,4$ coordinate directions, combined with optimal-length condenser lists, give deterministic recovery lists of only $315,612,1228$ square systems.  Adding fixed basis-first signing lists gives universal-forgery failure bounds below $2^{-153.502}$, $2^{-216.298}$, and $2^{-279.095}$.  The full coordinate and pair-line lists make those bounds much smaller.

The projective route remains the fastest proved asymptotic route.  It reduces recovery to a smooth complete intersection of degree $2^V$ and uses a finite-field Kronecker solver with expected bit complexity arbitrarily close to quadratic in that degree.  Correctly coupling the solver exponent to its required working field gives a rigorous checked single-invocation forger with success at least
\[
0.999999996259,\qquad0.999999999998,\qquad0.999999964716
\]
in the public-random model at Levels I, III, and V.  The older straight-line-program theorem and the independent locally closed-set theorem support the same growth trend without being used to borrow unstated failure constants.

The normalized asymptotic indicators lie below the Round-2 category targets at all three levels.  They are not concrete gate counts: hidden constants, field arithmetic, and memory traffic still require separate evaluation.  The evidence is strongest at Levels III and V, where several independent theorem interfaces leave substantial exponent margins.  Level I is more sensitive to low-level costs, although the primary modern-solver indicator and the epsilon-free arithmetic cross-check are both below its target.

Space is the main obstacle to implementation.  A conventional degree-$2^V$ quotient representation is astronomically large even though the recovered key is compact.  The norm-gradient and resultant-gradient formulas isolate a much smaller output target, but no current method constructs that output without first building a huge algebraic state.  A low-space solver or direct root-extraction method would therefore change the practical picture more than a small improvement in the displayed time exponent.

\appendix

\section{Official seeded expansion: ideal primitives and fixed functions}\label{app:seeded-transfer}

The key-by-key theorem needs no expansion model.  This appendix asks how often official seeded key generation satisfies the sufficient algebraic conditions for deterministic equivalent-key recovery.  The official expansion uses public counter $2i-1$ for the symmetric block $A_i$, public counter $2i$ for the rectangular block $E_i$, and secret counter $1$ for the graph $G$~\cite[Secs.~4.7--4.10]{QRUOVSpec}.  For every recommended set, the secret buffer length $\tau_2$ is at most the public $A_i$ buffer length $\tau_1$.  Hence, if the independently sampled public and secret seeds happen to coincide, the secret counter-$1$ raw stream is a prefix of the public $A_1$ stream.  The resulting correlation between $G$ and $A_1$ is expressly allowed by \cref{def:ideal}; it is not a bad event.

For Levels I, III, and V, let $\lambda\in\{128,192,256\}$ be the seed length.  The specification uses $m$ public buffers of length $\tau_1$ bytes and $m$ public buffers of length $\tau_2$ bytes.  Each public rejection sampler exhausts its prescribed buffer with probability at most $2^{-\lambda}$.  Put
\begin{equation}\label{eq:public-rejection}
\varepsilon_{\rm pub}:=2m2^{-\lambda}.
\end{equation}
No secret-buffer exhaustion term is needed: \cref{def:ideal} permits an arbitrary conditional law for $G$.

\subsection{SHAKE in the ideal-XOF model}

\begin{proposition}[Exact public-random transfer in the ideal-XOF model]\label{prop:rom-transfer}
Model SHAKE as a random function from each domain-separated input to one consistent infinite output string.  The total-variation distance between the official expanded-key distribution and the public-random model of \cref{def:ideal} is at most
\[
\delta_{\rm XOF}:=\varepsilon_{\rm pub}=2m2^{-\lambda}.
\]
This statement includes keys with equal public and secret seeds.
\end{proposition}

\begin{proof}
The $2m$ public seed/counter inputs are pairwise distinct, so their ideal-XOF strings are independent and uniform.  If the seeds differ, the secret string is independent of all public strings.  If they coincide, the secret counter-$1$ string is the required prefix of the public $A_1$ string; after rejection sampling, this induces an arbitrary conditional distribution of $G$ given $A_1$, exactly as permitted by \cref{def:ideal}.  In either case, the even-counter $E_i$ strings are independent uniform strings conditioned on $(A,G)$.  Abort only if one of the $2m$ public buffers contains too few accepted values.  The union bound \eqref{eq:public-rejection} pays for that event.  Conditioned on no public exhaustion, the accepted $A_i$ values have their product-uniform marginal and the accepted $E_i$ values are product-uniform conditioned on $(A,G)$, so the distribution is exactly \cref{def:ideal}.
\end{proof}

\subsection{AES in the ideal-cipher model}

Let
\[
T:=m\left(\left\lceil\frac{\tau_1}{16}\right\rceil
          +\left\lceil\frac{\tau_2}{16}\right\rceil\right)
\]
be the number of $128$-bit public cipher blocks, and define
\begin{equation}\label{eq:ideal-cipher-distance}
\delta_{\rm perm}(T)
:=1-\prod_{j=0}^{T-1}\left(1-\frac{j}{2^{128}}\right),
\qquad
\delta_{\rm IC}:=\delta_{\rm perm}(T)+\varepsilon_{\rm pub}.
\end{equation}

\begin{proposition}[Exact public-random transfer in the ideal-cipher model]\label{prop:icm-transfer}
Model AES as an ideal cipher with an independent random permutation for every key.  The total-variation distance between the official expanded-key distribution and the public-random model of \cref{def:ideal} is at most $\delta_{\rm IC}$.  This statement also includes equal public and secret AES keys.
\end{proposition}

\begin{proof}
For every fixed public key, the specification's $T$ public counter blocks are distinct.  Their ideal-cipher outputs are therefore a uniform ordered sample without replacement from $2^{128}$ blocks.  The exact total-variation distance from $T$ independent uniform blocks is the collision probability of the independent sample, namely $\delta_{\rm perm}(T)$.

Couple the without-replacement public block vector to an independent-block vector at this distance.  If the public and secret keys differ, generate $G$ from the independent secret-key permutation in both experiments.  If the keys coincide, generate $G$ from the prescribed prefix of the coupled public $A_1$ block vector in each experiment.  In both cases this is a common randomized postprocessing kernel, so total variation cannot increase.  Under the independent-block vector, $A$ has the product-uniform marginal and $E$ is product-uniform conditioned on $(A,G)$.  Truncation and rejection sampling are further postprocessing.  Finally, public-buffer exhaustion contributes at most \eqref{eq:public-rejection}.
\end{proof}

\begin{table}[t]
\centering
\small
\setlength{\tabcolsep}{4.3pt}
\caption{Official public-buffer parameters and exact ideal-primitive transfer bounds.  The last column gives base-two logarithms.  No seed-collision term is present.}
\label{tab:seeded-model-bounds}
\begin{tabular}{@{}c r r r r r r@{}}
\toprule
Level & $\lambda$ & $m$ & $\tau_1$ & $\tau_2$ & $T$ & $\log_2\delta_{\rm XOF}\ \big/\ \log_2\delta_{\rm IC}$\\
\midrule
I   &128&54 &4267 &2916 &24300  &$-121.245\ /\ -99.863$\\
III &192&78 &9020 &6123 &73866  &$-184.715\ /\ -96.655$\\
V   &256&105&16144&11018&178290 &$-248.286\ /\ -94.112$\\
\bottomrule
\end{tabular}
\end{table}

\subsection{From key-distribution distance to verified recovery}

Let
\[
 \eta_{\rm fast}:=\eta_{\off}+\eta_{\oil}+\eta_{\rm span}+\eta_{\rm sign},
 \qquad
 \eta_{\rm local}:=\tau_{m,V}(Q)^O+\eta_{\rm sign}
\]
be the public-random exceptional-key probabilities for the fast projective and graph-direction equivalent-key theorems.  Total variation is applied only to these key-level events, not to a bounded attack execution.

\begin{corollary}[Official seeded recovery in ideal-primitive models]\label{cor:seeded-ideal-primitives}
Under the official seeded key generation:
\begin{enumerate}
\item with SHAKE modeled as an ideal XOF, the probability that a key fails the graph-direction equivalent-key hypotheses is at most $\eta_{\rm local}+\delta_{\rm XOF}$, while the corresponding fast-projective bound is $\eta_{\rm fast}+\delta_{\rm XOF}$;
\item with AES modeled as an ideal cipher, the corresponding bounds are $\eta_{\rm local}+\delta_{\rm IC}$ and $\eta_{\rm fast}+\delta_{\rm IC}$.
\end{enumerate}
Every other key for the relevant route satisfies the corresponding key-by-key recovery theorem: accepted outputs are always correct, and verified restart terminates almost surely with the expected-work bounds of \cref{thm:graph-attack,thm:attack}.  At Levels I, III, and V, both ideal-key exceptions are negligible compared with the transfer terms, so the base-two logarithms round respectively to
\[
(-121.245,-184.715,-248.286)
\quad\text{and}\quad
(-99.863,-96.655,-94.112)
\]
for SHAKE and AES.
\end{corollary}

\begin{proof}
By \cref{prop:rom-transfer,prop:icm-transfer}, the probability of every key event changes by at most the corresponding total-variation distance.  Apply this separately to the graph-direction event of \cref{cor:graph-ideal-recovery} and the fast-projective event of \cref{cor:ideal-recovery}.  On either complement, the corresponding theorem is key by key and independent of how the expanded key was sampled.  Exact verification is distribution-free.  The ideal exceptional probabilities are far below the transfer terms at the displayed precision.
\end{proof}

\subsection{Why this does not prove a fixed-primitive theorem}

\begin{proposition}[Support barrier for deterministic seeded expansion]\label{prop:seed-support-barrier}
Let $E:\{0,1\}^{\lambda}\to\Omega$ be any deterministic public expansion map, let $S$ be uniform on $\{0,1\}^{\lambda}$, and let $U$ be uniform on the finite set $\Omega$.  Then
\[
\Delta_{\rm TV}(E(S),U)
\ge 1-\frac{2^{\lambda}}{|\Omega|}.
\]
In particular, if $\Omega=\{0,1\}^B$, the distance is at least $1-2^{\lambda-B}$.
\end{proposition}

\begin{proof}
The distribution $E(S)$ is supported on at most $2^\lambda$ points.  Let $R$ be its support.  Then $\Pr[E(S)\in R]=1$, while $\Pr[U\in R]=|R|/|\Omega|\le2^\lambda/|\Omega|$.  The variational characterization of total variation gives the claim.
\end{proof}

Thus the fixed public SHAKE/AES stream vector cannot be information-theoretically close to independent uniform streams.  Nor is there a generic public-seed PRG reduction: the public seed is revealed, expansion is deterministic, and a distinguisher can recompute the official stream and compare it with the transcript.  The ideal-XOF and ideal-cipher key-density conclusions above are exact statements inside those models; they do not instantiate themselves for fixed primitives.

The key-by-key recovery theorem nevertheless applies to every fixed-primitive key satisfying the key-by-key conditions.  A high-probability standard-model statement must therefore analyze the relevant algebraic bad events under the actual fixed expansion map, or introduce an explicitly event-specific hardness assumption.  The support barrier does not rule out such work; it rules out a generic statistical or transcript-pseudorandomness transfer that replaces the revealed deterministic expansion by independent uniform streams.

\section{Optional identification with the planted oil space}\label{app:uniqueness}

Neither recovery theorem requires uniqueness: it keeps every verified candidate and searches all of them for a signing witness.  This appendix proves the stronger generic identification statement that, for almost every key, no second $K$-rational common isotropic $O$-space exists.  Fix the planted oil space $\cO_0$.  Let $\cO'\ne\cO_0$ be another $O$-space and put
\[
s=O-\dim(\cO'\cap\cO_0),\qquad 1\le s\le O.
\]

\begin{lemma}[Relative Grassmannian stratum]\label{lem:relative-grassmann}
The locus of such $\cO'$ with fixed $s$ has dimension $s(V+O-s)$.
\end{lemma}

\begin{proof}
Choose the $(O-s)$-dimensional intersection inside $\cO_0$, the $s$-dimensional image of $\cO'$ in $K^{V+O}/\cO_0\simeq K^V$, and the graph map back to the $s$-dimensional quotient of $\cO_0$.  The dimensions are $(O-s)s$, $s(V-s)$, and $s^2$, respectively.
\end{proof}

\begin{lemma}[Independent extra isotropy conditions]\label{lem:extra-isotropy}
For fixed $\cO'$ in this stratum, requiring one central form that already vanishes on $\cO_0\times\cO_0$ also to vanish on $\cO'\times\cO'$ imposes exactly
\[
c_s=\binom{O+1}{2}-\binom{O-s+1}{2}=\frac{s(2O-s+1)}2
\]
additional independent linear conditions.
\end{lemma}

\begin{proof}
Restriction to $\cO'$ already vanishes on $\Sym^2(\cO'\cap\cO_0)$.  Conversely, after choosing complements, every symmetric form on $\cO'$ that vanishes on this subspace extends to an ambient symmetric form vanishing on $\cO_0\times\cO_0$.  The restriction map therefore has the displayed rank.
\end{proof}

\begin{theorem}[Uniqueness codimension and rational first moment]\label{thm:uniqueness}
For $m$ independent central forms, the key locus admitting an alternative common isotropic $O$-space of type $s$ has codimension at least
\[
\delta_s=m\frac{s(2O-s+1)}2-s(V+O-s).
\]
The minimum over $s$ equals $903,1927,3539$ for Levels~I, III, and V.  With Gaussian binomial coefficients denoted by $\genfrac{[}{]}{0pt}{}{a}{b}_Q$,
\[
\Pr[\text{a second $K$-rational oil space}]
\le
\sum_{s=1}^{O}
\genfrac{[}{]}{0pt}{}{O}{s}_Q
\genfrac{[}{]}{0pt}{}{V}{s}_Q
Q^{s^2-mc_s}.
\]
Writing the displayed first-moment bound as $\eta_{\rm uniq}$, we have
\[
\log_2\eta_{\rm uniq}<-18932.346,\quad -40401.586,\quad -74198.865
\]
at Levels I, III, and V, respectively.
\end{theorem}

\begin{proof}
Over the stratum of dimension $s(V+O-s)$, each of the $m$ form fibers loses $c_s$ dimensions by \cref{lem:extra-isotropy}.  The incidence variety therefore has codimension at least $mc_s-s(V+O-s)$ after projection to key space.  The function $\delta_s$ is concave in $s$, so its minimum on $1\le s\le O$ occurs at an endpoint; direct evaluation gives the displayed values and shows that $s=1$ minimizes in all three parameter sets.

For the finite-field statement, the number of $K$-rational $O$-spaces in the type-$s$ stratum is
\[
\genfrac{[}{]}{0pt}{}{O}{s}_Q
\genfrac{[}{]}{0pt}{}{V}{s}_Q Q^{s^2}.
\]
For each fixed $\cO'$, the $m$ independent forms satisfy the $c_s$ extra conditions with probability $Q^{-mc_s}$.  The union bound gives the formula; exact evaluation gives the stated logarithms.
\end{proof}

\begin{corollary}[Generic identification with the planted decomposition]\label{cor:planted-identification}
In the public-random model, outside the exceptional event of \cref{cor:graph-ideal-recovery} and the additional event of probability $\eta_{\rm uniq}$ in \cref{thm:uniqueness}, every exact graph-direction candidate has the planted oil space.  The same conclusion holds for the fast projective route outside the exceptional event of \cref{cor:ideal-recovery} and the uniqueness event.  In either case, the returned decomposition is the planted expanded decomposition up to the verified coordinate representation.  The secret seed is still not recovered or needed.
\end{corollary}

\begin{proof}
Every exact candidate defines a common isotropic $O$-space.  Uniqueness forces that space to be $\cO_0$, and its graph determines the expanded decomposition used by \cref{prop:verified-key}.
\end{proof}

\section{Alternative elimination and low-space directions}\label{app:alternatives}

The main text uses complete univariate resolutions because they are the only route for which we have both a suitable theorem and full failure accounting.  This appendix records two exact ways to extract a root more compactly once a quotient algebra exists, and then explains why several natural attempts do not yet improve the proved exponent.

\subsection{Extracting an isolated root from an existing quotient algebra}\label{subsec:implicit-extraction}

Let $A=L[y_1,\ldots,y_V]/I$ be the reduced affine coordinate algebra of a square core, of degree $d\le D$.  Write $Z$ for its geometric root set, and let $q_1,\ldots,q_m$ denote the original restricted equations in $A$.

\begin{proposition}[Norm-gradient root formula]\label{prop:implicit-root}
Suppose $p\in L^V$ is the unique point of $Z$ satisfying all $m$ original equations.  Choose $\rho\leftarrow L^m$, put $h=\sum_i\rho_iq_i\in A$, and define
\begin{equation}\label{eq:norm-jet}
\mathcal N_h(z_0,\ldots,z_V)
:=\operatorname{Norm}_{A/L}\!\left(h+z_0+\sum_{j=1}^Vz_jy_j\right).
\end{equation}
With probability at least $1-(d-1)/|L|$, $h$ vanishes at $p$ and at no other geometric core point.  On that event,
\begin{equation}\label{eq:norm-gradient-root}
\partial_{z_0}\mathcal N_h(0)\ne0,
\qquad
p_j=
\frac{\partial_{z_j}\mathcal N_h(0)}
     {\partial_{z_0}\mathcal N_h(0)}
\quad(1\le j\le V).
\end{equation}
Equivalently, if $M_h$ and $M_{y_j}$ are the multiplication operators in $A$, then $\ker M_h$ is one-dimensional and every nonzero $w\in\ker M_h$ satisfies $M_{y_j}w=p_jw$.
\end{proposition}

\begin{proof}
After scalar extension to an algebraic closure, reducedness gives
\[
A\otimes_L\overline L\cong\prod_{x\in Z}\overline L.
\]
At every $x\ne p$, the vector $(q_1(x),\ldots,q_m(x))$ is nonzero.  The condition $h(x)=0$ cuts out a proper $L$-linear subspace of $L^m$, so a uniform $\rho$ satisfies it with probability at most $1/|L|$.  A union bound over at most $d-1$ points proves the selector bound.  On the good event, the norm factors as
\[
\mathcal N_h(z)=
\prod_{x\in Z}\left(h(x)+z_0+\sum_{j=1}^Vz_jx_j\right).
\]
Differentiating at the origin gives
\[
\partial_{z_0}\mathcal N_h(0)=\prod_{x\ne p}h(x),
\qquad
\partial_{z_j}\mathcal N_h(0)=p_j\prod_{x\ne p}h(x),
\]
which proves \eqref{eq:norm-gradient-root}.  In the same product decomposition, $M_h$ is diagonal with a unique zero eigenvalue and the coordinate operators act on its kernel by the scalars $p_j$.
\end{proof}

\begin{remark}[Several surviving roots]\label{rem:norm-multiple}
For an arbitrary selector $h$, let $S=\{x\in Z:h(x)=0\}$.  The lowest nonzero homogeneous part of \eqref{eq:norm-jet} at the origin is
\[
\left(\prod_{x\notin S}h(x)\right)
\prod_{x\in S}\left(z_0+\sum_{j=1}^Vz_jx_j\right).
\]
Thus the first nonzero norm jet is the Chow product of all selected roots.  The linear formula \eqref{eq:norm-gradient-root} is exactly the case $|S|=1$.
\end{remark}

\begin{proposition}[Resultant-gradient root formula]\label{prop:resultant-gradient}
Assume $\operatorname{char}L\ne2$.  Let $F_0,\ldots,F_V\in L[t_0,\ldots,t_V]$ be homogeneous quadrics with a unique geometric common point $p=[p_0:\cdots:p_V]$.  Assume that their value Jacobian on the tangent space $T_p\PP^V$ has rank $V$.  Let $\mathcal R$ be their multivariate resultant, written in the monomial coefficients $c_{i,\alpha}$ of $F_i$, $|\alpha|=2$.  Then the coefficient tuple is a smooth point of the resultant hypersurface.  If $\lambda=(\lambda_0,\ldots,\lambda_V)$ spans the left kernel of the tangent Jacobian, there is a scalar $\kappa\ne0$ such that
\begin{equation}\label{eq:resultant-gradient}
\frac{\partial\mathcal R}{\partial c_{i,\alpha}}
=\kappa\lambda_i p^\alpha
\qquad(0\le i\le V,\ |\alpha|=2).
\end{equation}
In particular, for any $i,k$ with $\lambda_i p_k\ne0$,
\begin{equation}\label{eq:resultant-ratio}
\frac{p_j}{p_k}
=
\frac{\partial\mathcal R/\partial c_{i,e_j+e_k}}
     {\partial\mathcal R/\partial c_{i,2e_k}}
\qquad(0\le j\le V).
\end{equation}
The coefficient convention in \eqref{eq:resultant-ratio} treats $t_jt_k$ as one monomial, with no symmetric-matrix factor of two.
\end{proposition}

\begin{proof}
Consider the incidence variety of coefficient tuples and projective points satisfying $F_i(p)=0$ for every $i$.  Its first-order equations at the given pair are
\[
\delta F_i(p)+dF_i(p)[\delta p]=0\qquad(0\le i\le V).
\]
The tangent Jacobian has rank $V$, so a tangent displacement $\delta p$ exists exactly when
\[
\sum_{i=0}^V\lambda_i\,\delta F_i(p)=0.
\]
Uniqueness and transversality make the incidence projection locally one-to-one onto a smooth branch of the resultant hypersurface.  Its tangent hyperplane is therefore the kernel of the displayed nonzero functional.  The differential of the irreducible resultant is proportional to that functional, giving \eqref{eq:resultant-gradient}; the ratio formula follows by evaluating the two quadratic monomials at $p$.
\end{proof}

\begin{remark}[Compression achieved, evaluator still missing]\label{rem:implicit-memory}
The norm gradient outputs one root using only $V+1$ field elements, and the resultant gradient uses one nonzero coefficient block.  If quotient multiplication is available as a black box, a one-dimensional kernel can be found with $O(d)$ stored field elements and $O(d)$ multiplication-operator products by standard black-box linear algebra~\cite{Wiedemann}.  Alternatively, Baur--Strassen converts any arithmetic circuit for the norm or resultant into a circuit for all first derivatives with constant-factor overhead~\cite{BaurStrassen}.  These statements compress the \emph{root-extraction interface}; they do not construct the degree-$d$ quotient algebra, multiplication operators, norm, or resultant in $o(d)$ space.  With current constructions, the time remains near-quadratic in $d$ (with the fixed $\varepsilon$ loss of \cref{prop:finite-field-kronecker}) and at least one degree-$d$ state vector remains.  Moreover, \cref{prop:generic-full-isolation} gives no production-parameter density bound for uniqueness of the full-system root.  The formulas are therefore exact algorithmic targets, not a practical attack.
\end{remark}

\subsection{Why the spare equations do not yet lower the finite exponent}\label{subsec:exponent-audit}

The full restricted system has $e=m-V\in\{2,2,3\}$ equations beyond the square core, and generically those equations isolate the planted point.  We tested whether they also lower the solving exponent.  The following deliberately optimistic semi-regular screen uses the Hilbert series
\[
\frac{(1-t^2)^{V+e}}{(1-t)^{V+1}}
\]
for $V+e$ homogeneous quadrics in $V+1$ projective variables.  Because the actual system contains a planted point, this is a diagnostic rather than a theorem.  At the last positive coefficient, the monomial counts are:
\begin{center}
\small
\begin{tabular}{@{}c r r r r r@{}}
\toprule
Level & $V$ & degree & $\log_2 N$ & optimistic $\log_2N^2$ & current $\log_2D^2$\\
\midrule
I   &52  &27 &$69.787$  &$139.574$ &$104$\\
III &76  &39 &$102.585$ &$205.170$ &$152$\\
V   &102 &48 &$131.814$ &$263.629$ &$204$\\
\bottomrule
\end{tabular}
\end{center}
Even charging only $N^2$ sparse black-box work is worse than the baseline $D^2$ geometric scale and also worse than the primary $\varepsilon=0.3$ fast-line indicators after field and circuit conversion; dense linear algebra is worse still.  This does not rule out a coefficient-specific method, but it rules out the most direct overdetermined-Macaulay improvement under standard linear algebra.

\begin{proposition}[Underslicing loses more in hit probability than it saves in degree]\label{prop:underslice}
Let $0\le d<V$, and choose a uniform projective $d$-plane over $K=\F_Q$.  Its probability of meeting the planted projective oil space is
\begin{equation}\label{eq:underslice-probability}
\pi_d
=1-\prod_{i=0}^{d}
\frac{1-Q^{i-V}}{1-Q^{i-(V+O)}}.
\end{equation}
If the same degree-squared square-core solver is run independently on every sampled plane, its expected leading degree factor is $4^d/\pi_d$.  For $d=V-r$ and the \QRUOV{} base field $Q=127^3$, removing $r=1,2,3$ slice dimensions increases the base-two exponent relative to the $d=V$ slice by
\[
18.966053,\qquad37.932107,\qquad56.898161
\]
bits, respectively, before any other work is charged.  Choosing the slices over the larger solver field only worsens this tradeoff.
\end{proposition}

\begin{proof}
Write $k=d+1$.  A vector $k$-subspace disjoint from the fixed $O$-dimensional oil space is the graph of a linear map from a $k$-subspace of the $V$-dimensional quotient to the oil space.  Hence there are
\[
Q^{Ok}\genfrac{[}{]}{0pt}{}{V}{k}_Q
\]
disjoint subspaces, out of $\genfrac{[}{]}{0pt}{}{V+O}{k}_Q$ total.  Taking the ratio and expanding the Gaussian binomial coefficients gives \eqref{eq:underslice-probability}.  Independent repetition needs expected $1/\pi_d$ planes.  Substitution into $4^d/\pi_d$ gives the displayed values, which are recomputed by the companion numerical audit.  Asymptotically each removed dimension costs $\log_2Q-2=18.966054\ldots$ bits: it saves two degree-squared bits but loses roughly one factor of $Q$ in intersection probability.
\end{proof}

The same calculation covers the equivalent hybrid strategy of guessing one base-field coordinate of the oil point before solving: every guessed coordinate costs roughly $\log_2Q$ bits and saves only two bits in the degree-squared solver.  Thus the minimal guaranteed slice is also the best random underslice under the current solver model.

The other natural routes also stop short of an exponent improvement.  Explicit resultant matrices remain exponentially larger than the $D$-dimensional quotient representation in the dense quadratic regime~\cite{DAndreaDickenstein}.  Chow-form algorithms are polynomial in geometric degree but do not provide a subquadratic dependence specialized to these parameters~\cite{JeronimoEtAl}.  Faster modular composition lowers an operation used after a univariate algebra exists, not the geometric-intersection stage that constructs that algebra~\cite{NeigerEtAl}.

We also checked two recent refinements of Macaulay/Gr\"obner solving.  Degree smoothing replaces a full Macaulay matrix by a sufficient submatrix lying between two adjacent integer degrees~\cite{DegreeSmoothing}.  Our optimistic screen already charges the complete monomial space at the \emph{lower} adjacent degree and is still $35.574$, $53.170$, and $59.629$ bits above the $D^2$ geometric scale; smoothing within that interval cannot reverse the comparison.  The recent F4-tracer analysis for generic square degree-$\delta$ systems has leading arithmetic scale $\widetilde O(\delta^{\omega V+1}e^{cV})$ and relies on trace compatibility and genericity conjectures~\cite{F4Tracer}.  For quadrics, even its factor $2^{\omega V}$ exceeds $D^2=2^{2V}$ for every currently proved $\omega>2$.  Thus these improvements are valuable in their intended regimes but do not lower the exponent here.  The norm and resultant gradients identify a smaller \emph{output} target; no proved route found in this audit computes that target below the published near-quadratic-in-$D$ line used here.

\Needspace{10\baselineskip}
\section{Reproducibility and executable attack description}\label{app:algorithms}

\subsection{Companion numerical audit}

The companion file \path{QRUOV_numerical_audit.py} uses only the Python standard library.  It recomputes, from the three parameter tuples:
\begin{itemize}
\item the field size, B\'ezout degrees, Wronskian-list lengths, deterministic system counts, and exact random-matrix rank probabilities;
\item the basis-first and pair-line signing bounds, including the compact fixed-list bounds;
\item the minimum and safety extension degrees for the finite solver, its explicit outer and internal failure terms, and the one-invocation success probabilities;
\item every normalized cost indicator, robustness threshold, seeded-model transfer exponent, and memory figure quoted in the paper.
\end{itemize}
The script contains frozen assertions for the headline values and prints a plain-text report.  It is an arithmetic audit, not an implementation of the full-dimensional Kronecker solver.

\subsection{Executable attack description}

\paragraph{Graph-direction route.}
\begin{enumerate}
\item Expand the public key to matrices $P_i$ over $K=\F_{127^3}$.  Choose either a random nonzero direction $c\in K^O$ or a direction from the fixed coordinate list.  Form $q_{i,c}(x)=P_i(-x,c)$.
\item Choose either a random Wronskian condenser value $t$ or one of the $V(m-V)+1$ fixed values from \cref{cor:deterministic-outer-list}.  Form the square core $f=\mathsf C(t)q_c$ and $h=\det J_f$.
\item Over a field satisfying \eqref{eq:local-field-condition}, run the locally closed Kronecker algorithm on $f_1,\ldots,f_V$ localized at $h$.  Reject malformed, incomplete, or inconsistent output; a failed solver call causes only a retry.
\item Re-evaluate all $m$ direction equations at every represented root with $h(x)\ne0$.  For each survivor, solve the $O$ public linear systems of \cref{prop:graph-direction-completion}.  Keep only graphs passing every graph identity, total-isotropy check, inverse pull-back, and exact public-map factorization.
\item Test the basis signing assignments first.  If necessary, scan the remaining nonzero points on the designated pair lines.  Return the first exact candidate with an invertible oil matrix, together with its inverse.  For a fully deterministic outer list, repeat the finitely many probabilistic solver calls in rounds.
\end{enumerate}
The compact public-random version uses the first $3,4,4$ coordinate directions and the first $22,31,40$ basis signing assignments.  The full deterministic fallback uses all $O$ coordinate directions and all designated pair lines.

\paragraph{Fast projective route.}
\begin{enumerate}
\item Set $\varepsilon=0.3$.  Use the minimum fields $s=8,12,15$ for expected-restart analysis or the safety fields $s=9,13,16$ for the one-invocation theorem.  Sample a full-rank slice matrix $M$ and an output recombination $R$ over $L_s$; restrict the public quadrics and form the square core.
\item Run one finite checked invocation of the van der Hoeven--Lecerf solver.  Return \textsf{fail} on an exceptional division, degree mismatch, or consistency failure.  Do not hide an unbounded internal restart inside one outer draw.
\item Canonicalize and verify the resolution.  Evaluate all original restricted equations in its quotient algebra, prefilter with \cref{prop:separator-prefilter}, extract $L_s$-rational roots, and recheck all equations at each reconstructed point.
\item At every survivor, compute the common derivative kernel.  Require dimension $O$, descent to a graph over $K$, total isotropy, all graph identities, the official inverse pull-back, and exact factorization of the original public map.
\item Run the basis-first signing scan and, if needed, the pair-line fallback.  Return the first verified signing certificate; otherwise begin a fresh outer trial.
\end{enumerate}

Both wrappers have one-sided error: a solver failure may cause a retry, but exact public verification prevents an invalid equivalent key or invalid signature from being returned.  The output is a verified expanded signing decomposition and an inverse oil matrix, not the original seed or an official serialized private key.

\begin{multicols}{2}
\raggedcolumns
\begin{thebibliography}{99}
\fontsize{8.5pt}{9.6pt}\selectfont
\setlength{\itemsep}{0pt}
\setlength{\parsep}{0pt}

\bibitem{QRUOVSpec}
H.~Furue, Y.~Ikematsu, F.~Hoshino, T.~Takagi, H.~Kosuge, K.~Yamakoshi, R.~Akiyama, S.~Nakamura, S.~Orihara, and K.~Kinjo.
\newblock \emph{QR-UOV Specification Document (Round 2 submission)}.
\newblock NIST Additional Digital Signatures Project, 5 February 2025.

\bibitem{FurueEtAl2026}
H.~Furue, Y.~Ikematsu, F.~Hoshino, T.~Takagi, H.~Kosuge, K.~Yamakoshi, R.~Akiyama, S.~Nakamura, S.~Orihara, and K.~Kinjo.
\newblock Post-quantum multivariate signature scheme QR-UOV: Security, parameters, and implementations.
\newblock \emph{Japan Journal of Industrial and Applied Mathematics}, 2026.
\newblock doi:10.1007/s13160-026-00809-7.

\bibitem{KPG1999}
A.~Kipnis, J.~Patarin, and L.~Goubin.
\newblock Unbalanced Oil and Vinegar signature schemes.
\newblock In \emph{EUROCRYPT 1999}, pp.~206--222, 1999.

\bibitem{QRUOV2021}
H.~Furue, Y.~Ikematsu, Y.~Kiyomura, and T.~Takagi.
\newblock A new variant of unbalanced oil and vinegar using quotient ring: QR-UOV.
\newblock In \emph{ASIACRYPT 2021, Part IV}, pp.~187--217, 2021.

\bibitem{NISTRound3}
National Institute of Standards and Technology.
\newblock Nine candidates advance to the third round of the Additional Digital Signatures for the PQC standardization process.
\newblock NIST CSRC, 14 May 2026.

\bibitem{NISTIR8610}
G.~Alagic et al.
\newblock Status report on the second round of the Additional Digital Signature Schemes for the NIST Post-Quantum Cryptography Standardization Process.
\newblock NIST IR~8610, May 2026.

\bibitem{Benoist}
O.~Benoist.
\newblock Degr\'es d'homog\'en\'eit\'e de l'ensemble des intersections compl\`etes singuli\`eres.
\newblock \emph{Annales de l'Institut Fourier}, 62(3):1189--1214, 2012.

\bibitem{vdHL}
J.~van der Hoeven and G.~Lecerf.
\newblock On the complexity exponent of polynomial system solving.
\newblock \emph{Foundations of Computational Mathematics}, 21(1):1--57, 2021.

\bibitem{GiustiLecerfSalvy2001}
M.~Giusti, G.~Lecerf, and B.~Salvy.
\newblock A Gr\"obner free alternative for polynomial system solving.
\newblock \emph{Journal of Complexity}, 17(1):154--211, 2001.
\newblock doi:10.1006/jcom.2000.0571.

\bibitem{Gimenez2026}
N.~Gim\'enez, J.~Heintz, G.~Matera, L.~M.~Pardo, M.~P\'erez, and M.~Privitelli.
\newblock A Kronecker algorithm for locally closed sets over a perfect field.
\newblock arXiv:2512.14888v2, 5 June 2026.

\bibitem{Beullens}
W.~Beullens.
\newblock Improved cryptanalysis of UOV and Rainbow.
\newblock In \emph{EUROCRYPT 2021, Part I}, pp.~348--373, 2021.

\bibitem{FurueIkematsu}
H.~Furue and Y.~Ikematsu.
\newblock A new security analysis against MAYO and QR-UOV using rectangular MinRank attack.
\newblock In \emph{IWSEC 2023}, pp.~101--116, 2023.

\bibitem{Suzuki}
T.~Suzuki, H.~Furue, T.~Ito, S.~Nakamura, and S.~Uchiyama.
\newblock An extended rectangular MinRank attack against UOV and its variants.
\newblock In \emph{IWSEC 2025}, LNCS~16208, pp.~111--130, Springer, 2026.

\bibitem{FurueIkematsu2026}
H.~Furue and Y.~Ikematsu.
\newblock Key recovery attacks on UOV using $p^\ell$-truncated polynomial rings.
\newblock Cryptology ePrint Archive, Paper 2026/298, 2026.

\bibitem{JustGuess}
A.~May, M.~Ostuzzi, and H.~Ressler.
\newblock Just Guess: Improved (quantum) algorithm for the underdetermined MQ problem.
\newblock In \emph{EUROCRYPT 2026, Part IV}, LNCS~16544, pp.~334--362, 2026.

\bibitem{QRUOVJustGuess}
S.~Orihara, on behalf of the \QRUOV{} team.
\newblock Official comment: \emph{Just Guess} estimates for the main $q=127$ parameter sets.
\newblock NIST PQC Forum, 20 January 2026.

\bibitem{Asanuma2026}
H.~Asanuma, Y.~Chen, H.~Furue, K.~Sakata, and T.~Takagi.
\newblock Improved complexity estimates for underdetermined MQ systems via generalized variable partitioning.
\newblock Cryptology ePrint Archive, Paper 2026/1054, 2026.

\bibitem{Pebereau}
P.~P\'ebereau.
\newblock One vector to rule them all: Key recovery from one vector in UOV schemes.
\newblock In \emph{PQCrypto 2024}, pp.~92--108, 2024; ePrint 2023/1131.


\bibitem{PippengerFischer}
N.~Pippenger and M.~J.~Fischer.
\newblock Relations among complexity measures.
\newblock \emph{Journal of the ACM}, 26(2):361--381, 1979.

\bibitem{BaurStrassen}
W.~Baur and V.~Strassen.
\newblock The complexity of partial derivatives.
\newblock \emph{Theoretical Computer Science}, 22(3):317--330, 1983.
\newblock doi:10.1016/0304-3975(83)90110-X.

\bibitem{DAndreaDickenstein}
C.~D'Andrea and A.~Dickenstein.
\newblock Explicit formulas for the multivariate resultant.
\newblock \emph{Journal of Pure and Applied Algebra}, 164(1--2):59--86, 2001.
\newblock doi:10.1016/S0022-4049(00)00145-6.

\bibitem{JeronimoEtAl}
G.~Jeronimo, T.~Krick, J.~Sabia, and M.~Sombra.
\newblock The computational complexity of the Chow form.
\newblock \emph{Foundations of Computational Mathematics}, 4(1):41--117, 2004.
\newblock doi:10.1007/s10208-002-0078-2.

\bibitem{NeigerEtAl}
V.~Neiger, B.~Salvy, E.~Schost, and G.~Villard.
\newblock Faster modular composition using two relation matrices.
\newblock In \emph{ISSAC 2026}, pp.~315--324, 2026; arXiv:2601.17422.

\bibitem{DegreeSmoothing}
S.~Jaques, L.~Ran, S.~Samardjiska, and M.~Seitner.
\newblock Smoothing the degree of regularity for polynomial systems.
\newblock Cryptology ePrint Archive, Paper 2026/408, 2026.

\bibitem{F4Tracer}
R.~Kouba, V.~Neiger, and M.~Safey El Din.
\newblock A complexity analysis of the F4 Gr\"obner basis algorithm with tracer data.
\newblock arXiv:2603.16378, 2026.

\bibitem{KaltofenShoup}
E.~Kaltofen and V.~Shoup.
\newblock Subquadratic-time factoring of polynomials over finite fields.
\newblock \emph{Mathematics of Computation}, 67(223):1179--1197, 1998.

\bibitem{Wiedemann}
D.~Wiedemann.
\newblock Solving sparse linear equations over finite fields.
\newblock \emph{IEEE Transactions on Information Theory}, 32(1):54--62, 1986.

\bibitem{Heintz}
J.~Heintz.
\newblock Definability and fast quantifier elimination in algebraically closed fields.
\newblock \emph{Theoretical Computer Science}, 24:239--277, 1983.

\bibitem{LachaudRolland}
G.~Lachaud and R.~Rolland.
\newblock On the number of points of algebraic sets over finite fields.
\newblock \emph{Journal of Pure and Applied Algebra}, 219(11):5117--5136, 2015.

\end{thebibliography}
\end{multicols}

\end{document}
